\documentclass[10pt,a4paper]{article}
\usepackage[utf8]{inputenc}
\usepackage[T1]{fontenc}
\usepackage[french]{babel}
\usepackage{lmodern}
\usepackage[margin=1.9cm]{geometry}
\usepackage{amsmath,amssymb}
\usepackage{booktabs,array,tabularx,longtable,ltablex}
\keepXColumns
\usepackage{xcolor}
\usepackage{fancyvrb}
\usepackage{caption}
\usepackage{tikz}
\usetikzlibrary{arrows.meta,positioning,fit}
\usepackage[hidelinks]{hyperref}
\usepackage{textcomp}
\usepackage{siunitx}
\sisetup{locale=FR,detect-all}

\setlength{\parskip}{0.4em}
\setlength{\parindent}{0pt}
\fvset{fontsize=\small,frame=single,framesep=2mm,rulecolor=\color{black!30}}
\newcommand{\cmd}[1]{\texttt{#1}}
% nom de paramètre ou de variable du code, utilisable en texte et en mathématiques
\newcommand{\vp}[1]{\text{\texttt{\detokenize{#1}}}}
\newcommand{\evA}{\textcolor{green!45!black}{\textbf{A}}}
\newcommand{\evB}{\textcolor{green!45!black}{\textbf{B}}}
\newcommand{\evC}{\textcolor{orange!80!black}{\textbf{C}}}
\newcommand{\evD}{\textcolor{red!70!black}{\textbf{D}}}
\newcommand{\acc}[1]{\textsl{#1}}
\newcommand{\limite}[1]{\par\textbf{Limites et écarts connus.} #1}
\newcommand{\sources}[1]{\par\textbf{Base bibliographique.} #1}
\newcolumntype{L}{>{\raggedright\arraybackslash}X}
\newcolumntype{P}[1]{>{\raggedright\arraybackslash}p{#1}}
\renewcommand{\arraystretch}{1.15}
\newcommand{\dd}{\mathrm{d}}

\title{Guide du modèle VetTwin\\[0.3em]
\large Patient canin virtuel \emph{vetsim} : équations, valeurs, sources et fonctionnement, bloc par bloc}
\author{Projet VetTwin}
\date{18 novembre 2026}

\begin{document}
\sloppy
\maketitle

\begin{center}
\fbox{\parbox{0.92\linewidth}{\textbf{Usage pédagogique uniquement.} \emph{vetsim} est un patient canin virtuel destiné à l'enseignement. Ses paramètres sont en partie calés sur des moyennes publiées et ne sont pas validés cliniquement ; de nombreuses valeurs sont des \emph{placeholders} (niveau de preuve D). Aucun résultat du modèle ne doit servir à décider d'une dose, d'un traitement ou d'un pronostic réel.}}
\end{center}

\tableofcontents
\newpage

%=====================================================================
\section{Préambule}
%=====================================================================

\subsection{Objet et statut du document}
Ce guide décrit, bloc par bloc, les équations que le moteur résout, les valeurs des paramètres du chien sain et des modules, la base bibliographique de chaque choix et les limites connues. Les valeurs de sortie citées (FC, PAM, POG\ldots) proviennent des calages et des vérifications, avec leur date.

\subsection{Références normative et secondaire}
\begin{itemize}
\item \textbf{Référence normative} (décision du 18/11/2026, \cmd{docs/VERIFICATION\_MALADIES\_MEDICAMENTS.tex}, section « Gouvernance ») : le moteur Rust \cmd{crates/vetsim-core/src/}. Les équations de ce guide sont transcrites de ce code (\cmd{hemo.rs}, \cmd{engine.rs}, \cmd{gas.rs}, \cmd{autonomic.rs}, \cmd{rhythm.rs}, \cmd{pkpd.rs}, \cmd{kinetics.rs}, \cmd{signals.rs}, \cmd{pcg.rs}).
\item \textbf{Valeurs} : fichiers JSON de \cmd{modules/} (\cmd{healthy\_dog.json}, \cmd{diseases/*.json}, \cmd{drugs/*.json}, \cmd{auscultation.json}). Ils font foi pour les nombres.
\item \textbf{Référence secondaire} : le Python de \cmd{reference/python/}, conservé pour la contre-vérification (golden, fixtures). Il n'est cité ici que pour l'historique des calages (\cmd{fit\_hemo.py}, \cmd{fit\_gas.py}, \cmd{fit\_mmvd.py}, \cmd{fit\_pk.py}).
\end{itemize}

\subsection{Conventions d'unités}
Pression en mmHg ; volume en mL ; temps en s dans l'hémodynamique et la conduction, en h dans la PK ; débit en mL/s (le débit cardiaque est rapporté en L/min) ; concentrations en mmol/L, sauf albumine (g/dL en interne, g/L en sortie), créatinine (\textmu mol/L), urée (mmol/L), médicaments (mg/L) ; poids \vp{bw} en kg ; FC en /min ; élastance en mmHg/mL ; compliance en mL/mmHg ; résistance en mmHg·s/mL.

\textbf{Attention au symbole $s$.} Dans le patient sain, $s = \vp{bw}/20$ (facteur d'échelle allométrique, chien de référence de 20~kg). Dans un module de maladie, $s$ désigne la \emph{sévérité} (0 à 1). Les deux ne coexistent jamais dans une même formule.

\subsection{Niveaux de preuve et niveaux d'accès}
\textbf{Niveau de preuve} (champ \vp{evidence} des modules) :
\evA{} calé sur des mesures publiées ; \evB{} bornes ou critères publiés ; \evC{} plausibilité à partir d'intervalles publiés ; \evD{} placeholder pédagogique, sans source.

\textbf{Niveau d'accès} à la source, tel que les documents du dépôt le déclarent :
\acc{lu} (texte intégral lu, souvent dans PMC) ; \acc{résumé} (résumé seul) ; \acc{moteur} (extrait renvoyé par un moteur de recherche, page non ouverte) ; \acc{dépôt} (information tirée du code ou d'un module, sans source externe). Les DOI des entrées de \cmd{docs/journal/references.bib} ont été vérifiés dans OpenAlex le 17/11/2026 ; ceux qui apparaissent seulement dans les notes des modules ou dans les documents du 18/11 ne l'ont pas tous été (indiqué au cas par cas).

\textbf{Mention « non documenté ».} Quand ni le code, ni le module, ni un document du dépôt ne donne la source d'une valeur, ce guide l'écrit ; il n'invente ni source ni valeur.

%=====================================================================
\section{Vue d'ensemble}
%=====================================================================

\subsection{Blocs du modèle}
\cmd{blocks.rs} rattache chaque paramètre aux blocs (équations) où il intervient ; le test \cmd{every\_parameter\_belongs\_to\_a\_block} vérifie que tout paramètre a un bloc et que tout bloc a au moins un paramètre. Le tableau~\ref{tab:blocs} en donne la liste ; les sections suivantes suivent à peu près cet ordre.

\begin{tabularx}{\linewidth}{@{}P{3.6cm}LP{2.8cm}@{}}
\caption{Blocs du modèle (\cmd{blocks.rs}) et principaux paramètres.}\label{tab:blocs}\\
\toprule
Bloc (\cmd{Block}) & Paramètres principaux & Section\\
\midrule\endfirsthead
\toprule
Bloc (\cmd{Block}) & Paramètres principaux & Section\\
\midrule\endhead
Patient & \vp{bw} & \ref{sec:patient}\\
Ventricule gauche / droit & \vp{ees_lv}, \vp{v0_lv}, \vp{a_lv}, \vp{b_lv} ; \vp{ees_rv}, \vp{v0_rv}, \vp{a_rv}, \vp{b_rv} ; \vp{ts_coef}, \vp{ts_restitution}, \vp{ca_tau}, \vp{ca_recirc}, \vp{ees_wide} & \ref{sec:coeur}\\
Circulation systémique & \vp{c_ao}, \vp{c_sv}, \vp{ven_cap}, \vp{v0_ao}, \vp{v0_sv}, \vp{r_sys}, \vp{k_ra} & \ref{sec:circ}\\
Circulation pulmonaire et OG & \vp{c_pa}, \vp{c_pv}, \vp{v0_pa}, \vp{v0_pv}, \vp{r_pul}, \vp{k_la}, \vp{edema_k} & \ref{sec:circ}, \ref{sec:oreillettes}\\
Volémie & \vp{bvf}, \vp{vol_shift} & \ref{sec:circ}\\
Valves & \vp{r_valve}, \vp{g_mr}, \vp{g_ar}, \vp{g_tr}, \vp{g_pr}, \vp{mr_k}, \vp{mr_vref} & \ref{sec:valves}\\
Système autonome, baroréflexe et conduction & \vp{hr}, \vp{map_set}, \vp{baro_gain}, \vp{s0}, \vp{v0}, \vp{a_s}, \vp{a_v}, \vp{vagal_add}, \vp{vagal_block}, \vp{beta_block}, \vp{rsa_depth}, \vp{vag_noise}, \vp{sym_noise}, \vp{baro_sens} ; conduction : \vp{sa_block}, \vp{af}, \vp{av_erp}, \vp{av_symp}\ldots & \ref{sec:auto}, \ref{sec:rythme}\\
Ventilation et échanges gazeux & \vp{vco2}, \vp{va0}, \vp{vent_drive}, \vp{aagrad}, \vp{fio2}, \vp{pb}, \vp{rq}, \vp{k_weak}, \vp{k_weak_vent} & \ref{sec:gaz}\\
Transport de l'O$_2$ & \vp{hb}, \vp{p50}, \vp{hill_n} & \ref{sec:gaz}\\
Acide-base (Stewart) & \vp{na}, \vp{k}, \vp{cl}, \vp{xa}, \vp{alb}, \vp{phos} & \ref{sec:gaz}\\
Métabolisme et lactate & \vp{vo2}, \vp{lact0}, \vp{glucose}, \vp{hypox_pvo2}, \vp{hypox_ees} & \ref{sec:metab}\\
Rein & \vp{nephron_mass}, \vp{gfr0}, \vp{creat0}, \vp{urea0}, \vp{creat_vd}, \vp{urea_vd} & \ref{sec:rein}\\
Foie & \vp{hepatic_fn}, \vp{alt} & \ref{sec:foie}\\
ECG & \vp{pr}, \vp{qrs}, \vp{qt}, \vp{r_amp}, \vp{t_amp}, \vp{p_amp} (et K$^+$) & \ref{sec:signaux}\\
\bottomrule
\end{tabularx}

\subsection{Boucle d'évaluation}
La figure~\ref{fig:boucle} résume l'enchaînement d'une évaluation (\cmd{engine::evaluate}).

\begin{center}
\begin{tikzpicture}[
  box/.style={draw, rounded corners, align=center, minimum height=8mm, inner sep=3pt, font=\small, fill=black!4},
  arr/.style={-{Stealth[length=2mm]}, thick},
  node distance=6mm and 8mm]
\node[box] (pat) at (-5.2,0) {Patient sain\\\cmd{healthy\_dog.json}\\(allométrie, variabilité)};
\node[box] (dis) at (0,0) {Maladies\\(effets DSL\\sur les paramètres)};
\node[box] (pk) at (5.2,0) {PK/PD (\cmd{timecourse})\\$C(t)$, $D$, $U_\text{eff}$\\$\to$ \cmd{apply\_drug\_effects}};
\node[box, minimum width=4.2cm] (baro) at (0,-2.2) {Baroréflexe\\(\cmd{evaluate\_once}, $\le 5$ passages)};
\node[box] (cond) at (-2.8,-4.2) {Conduction\\(\cmd{rhythm.rs})};
\node[box] (hemo) at (2.8,-4.2) {Hémodynamique\\6 compartiments (\cmd{hemo.rs})};
\node[box] (met) at (0,-6.4) {Aval métabolique : gaz, Stewart,\\lactate, rein (\cmd{metabolic\_state})};
\node[box] (dec) at (-6.0,-6.4) {Décompensation\\(point fixe $f_c$, $f_r$)\\et viabilité};
\node[box] (sig) at (6.0,-6.4) {Signaux : VFC, ECG,\\PCG ; créatinine cinétique};
\draw[arr] (pat) -- (dis);
\draw[arr] (dis) -- (pk);
\draw[arr] (pk.south) |- (baro.east);
\draw[arr] (baro) -- node[left,font=\scriptsize]{FC, $s_\text{eff}$, $v_\text{eff}$} (cond);
\draw[arr] (cond) -- node[below,font=\scriptsize]{battements, PR} (hemo);
\draw[arr] (hemo) -- node[right,font=\scriptsize]{PAM $\to d$} (baro);
\draw[arr] (hemo) -- node[right,font=\scriptsize]{DC, PAM, POG} (met);
\draw[arr] (met) -- node[below,font=\scriptsize]{pH, PvO$_2$, PPC} (dec);
\draw[arr] (dec.north) |- node[pos=0.3,left,font=\scriptsize]{$E_{es}\cdot f_c$, $R_{sys}\cdot f_r$} (baro.west);
\draw[arr] (met) -- (sig);
\end{tikzpicture}
\captionof{figure}{Boucle d'évaluation. Les flèches de retour sont les deux points fixes : baroréflexe (commande $d$) et décompensation (facteurs $f_c$, $f_r$).}\label{fig:boucle}
\end{center}

\textbf{Étapes.}
\begin{enumerate}
\item \textbf{Patient.} Chaque paramètre de \cmd{healthy\_dog.json} est une expression de $s = \vp{bw}/20$ ; quatorze paramètres reçoivent un multiplicateur lognormal de moyenne 1 (section~\ref{sec:patient}).
\item \textbf{Maladies.} Chaque module tire une sévérité, ses tirages (\vp{draws}), calcule ses \vp{let}, puis applique ses effets dans l'ordre (\vp{mul} ou \vp{add}, \vp{when}, \vp{clamp}).
\item \textbf{Médicaments.} \cmd{pkpd::timecourse} intègre la PK et l'état de volume du diurétique ; \cmd{apply\_drug\_effects} applique, pour l'instant demandé, les effets thérapeutiques (variable $E$) puis toxiques (variable $T$).
\item \textbf{Évaluation} (\cmd{engine::evaluate}) : point fixe de décompensation autour de \cmd{evaluate\_once}, qui contient lui-même le point fixe du baroréflexe ; chaque passage du baroréflexe simule la conduction puis l'hémodynamique ; l'aval métabolique est calculé sur le dernier passage.
\item \textbf{Viabilité.} Un état non viable lève l'erreur \cmd{Error::NonViable} avec une cause (section~\ref{sec:decomp}).
\item \textbf{Signaux} (\cmd{signals.rs}, \cmd{autonomic.rs}, \cmd{pcg.rs}) : simulation autonome à 20~Hz (VFC), ECG construit depuis la conduction, sons du cœur ; créatinine et urée cinétiques dans les évolutions temporelles (\cmd{kinetics.rs}).
\end{enumerate}

\textbf{Intégration numérique.} Hémodynamique : Dormand--Prince 5(4) à pas adaptatif (\cmd{ode.rs}, interpolant et contrôle de pas de \cmd{scipy.integrate.RK45}), $\vp{rtol} = \vp{atol} = 10^{-6}$, pas maximal 4~ms ; PK : même intégrateur, $\vp{rtol} = 10^{-8}$, $\vp{atol} = 10^{-10}$, relancé à chaque dose et chaque instant d'enregistrement. Méthode explicite : un module rendant le système raide imposerait un autre solveur.

%=====================================================================
\section{Patient sain et allométrie}\label{sec:patient}
%=====================================================================

\textbf{Rôle.} Fournir les paramètres d'un chien sain de poids quelconque. \cmd{healthy\_dog.json} (version 1.0.0, \vp{reference_bw_kg} = 20, imposé par le schéma) donne chaque paramètre comme expression de $s$.

\textbf{Lois d'échelle retenues} (lecture des expressions du JSON) :
\begin{itemize}
\item FC, VO$_2$ et VCO$_2$ par kg : $\propto s^{-0{,}25}$ ($\vp{hr} = 87\,s^{-0{,}25}$, $\vp{vo2} = 5{,}5\,s^{-0{,}25}$~mL/min/kg, $\vp{vco2} = 4{,}6\,s^{-0{,}25}$) ;
\item volumes et compliances : $\propto s$ ; élastances et raideurs diastoliques ($\vp{ees}$, $b$) : $\propto 1/s$ ;
\item résistances : $\propto s^{-0{,}75}$ ($\vp{r_sys} = 1{,}7572/s^{0{,}75}$), si bien que les débits varient en $s^{0{,}75}$ ; la clairance des médicaments suit $\vp{bw}^{0{,}75}$ (section~\ref{sec:pk}) ;
\item le reste est indépendant du poids (concentrations, $a$ des ventricules, seuils).
\end{itemize}

\textbf{Variabilité interindividuelle.} Multiplicateur $z = \exp(\mathcal N(-\sigma^2/2, \sigma^2))$, de moyenne 1, appliqué à : \vp{ees_lv} ($\sigma$ 0,08), \vp{r_sys} (0,08), \vp{c_ao} (0,10), \vp{hr} (0,06), \vp{hb} (0,06), \vp{creat0} (0,10), \vp{gfr0} (0,08), \vp{aagrad} (0,20), \vp{vco2} (0,05), \vp{na} (0,006), \vp{k} (0,05), \vp{cl} (0,01), \vp{alb} (0,06), \vp{alt} (0,25). Source des $\sigma$ : non documentée (choix du module). Mode déterministe (golden, vérifications) : $z = 1$.

\textbf{Calage du chien sain.} Les valeurs ne sont pas toutes sourcées une à une. Deux scripts secondaires ont fixé des paramètres par moindres carrés sur des cibles publiées (\cmd{calibration/targets.json}) :
\begin{itemize}
\item \cmd{fit\_hemo.py} : \vp{r_sys}, \vp{vol_shift}, \vp{ees_lv}, \vp{r_pul}, \vp{b_lv} sur PAM 103, DC (IC 4,42 L/min/m$^2$ $\times$ surface corporelle), PAP 14, POG 5,5 (Haskins 2005) et FE 0,55 (source de la FE non documentée) ;
\item \cmd{fit\_gas.py} : \vp{xa} et \vp{va0} sur pH 7,38 et PaCO$_2$ 40,2 (Haskins 2005).
\end{itemize}
Les autres élastances, compliances et volumes non contraints n'ont pas de source dans le dépôt.

\begin{tabularx}{\linewidth}{@{}P{2.3cm}P{2.6cm}P{2.1cm}LP{0.6cm}@{}}
\caption{Paramètres du chien sain (\cmd{healthy\_dog.json}) : expression, valeur à 20~kg, source déclarée.}\label{tab:sain}\\
\toprule
Paramètre & Expression & Unité & Source (accès) & Év.\\
\midrule\endfirsthead
\toprule
Paramètre & Expression & Unité & Source (accès) & Év.\\
\midrule\endhead
\multicolumn{5}{@{}l}{\emph{Cœur}}\\
\vp{hr} & $87\,s^{-0{,}25}$ & /min & FC 87, Haskins 2005 (97 chiens vigiles ; \acc{résumé}) & A\\
\vp{ees_lv}, \vp{ees_rv} & $4{,}4915/s$ ; $1{,}6/s$ & mmHg/mL & \vp{ees_lv} calé (\cmd{fit\_hemo.py}) ; \vp{ees_rv} non documenté & C\\
\vp{v0_lv}, \vp{v0_rv} & $4\,s$ ; $6\,s$ & mL & non documenté & --\\
\vp{a_lv}, \vp{a_rv} & 1,0 ; 0,4 & mmHg & non documenté & --\\
\vp{b_lv}, \vp{b_rv} & $0{,}028137/s$ ; $0{,}05/s$ & 1/mL & \vp{b_lv} calé (\cmd{fit\_hemo.py}) & C\\
\vp{ts_coef}, \vp{act_gamma}, \vp{tr_frac} & 0,25 ; 0,55 ; 0,35 & -- & calés le 18/11/2026 sur les intervalles systoliques et diastoliques du chien vigile (tableau~\ref{tab:sti}) & C\\
\vp{emd_v} & 0,022 s & s & début de Q $\to$ début de la montée de pression VG $22 \pm 4$~ms, $n = 56$ (Kettunen 1985, \acc{résumé}) ; $\approx 30$~ms (Strik 2012, \acc{lu}) & B\\
\vp{emd_a} & 0 & s & aucune valeur chez le chien sain (PA-TDI seulement chez des cardiopathes, Neves 2018) & D\\
\vp{ts_restitution} & 0,1 & -- & non documenté & D\\
\vp{ca_tau} & 0,12 s & s & restitution mécanique in situ, thorax fermé, chien vigile instrumenté : $\tau = 122 \pm 26$~ms sous stimulation ventriculaire, $64 \pm 11$~ms sous stimulation atriale, $n = 9$ (Freeman \& Colston 1990, \acc{résumé}) ; 82~ms (Prabhu 1998) & C\\
\vp{ca_recirc} & 0,6 & -- & structure : Schouten 1987, Wier 1986 ; valeur : placeholder & D\\
\vp{ees_wide} & 0,8 & -- & activation hors Purkinje : d$P$/d$t_\text{max}$ $-22$~\% en BBG aigu (Vernooy 2007, cité par Strik 2012, \acc{lu}) ; pic de PVG $-26 \pm 10$~\% en stimulation de paroi libre VD (Burkhoff 1986, \acc{résumé}) ; PVG télésystolique $-11$~\% en stimulation de l'apex VD (Zhou 2014, \acc{lu}) & C\\
\multicolumn{5}{@{}l}{\emph{Vaisseaux et volémie}}\\
\vp{c_ao}, \vp{c_sv} & $0{,}7\,s$ ; $55\,s$ & mL/mmHg & non documenté & --\\
\vp{c_pa} & $2{,}5\,s$ & mL/mmHg & PAP du chien vigile 20,0/12,8/6,4~mmHg (systolique/moyenne/diastolique, médianes, $n = 7$ ; Yuchi 2021, \acc{lu}) ; modèle 19,5/13,7/7,8 (1,5 donnait 22/13,5/5,6 et un débit pulmonaire en diastole) & C\\
\vp{c_pv} & $5\,s$ & mL/mmHg & non documenté & --\\
\vp{v0_ao}, \vp{v0_sv}, \vp{v0_pa}, \vp{v0_pv} & $150\,s$ ; $900\,s$ ; $60\,s$ ; $120\,s$ & mL & non documenté & --\\
\vp{ven_cap} & 1,0 & -- & capacitance relative (1 = basale) & --\\
\vp{r_sys}, \vp{r_pul} & $1{,}7572/s^{0{,}75}$ ; $0{,}156/s^{0{,}75}$ & mmHg·s/mL & calés (\cmd{fit\_hemo.py}) sur Haskins 2005 & A\\
\vp{bvf}, \vp{vol_shift} & 1,0 ; $-21\,s$ & -- ; mL & \vp{vol_shift} calé (\cmd{fit\_hemo.py}) & C\\
\multicolumn{5}{@{}l}{\emph{Valves et oreillettes}}\\
\vp{r_valve} & $0{,}03/s^{0{,}75}$ & mmHg·s/mL & non documenté & --\\
\vp{g_mr}, \vp{g_ar}, \vp{g_tr}, \vp{g_pr} & 0,0005 ; 0,0002 ; 0,0002 ; 0,0002 & mL/s/mmHg & fuite résiduelle quasi nulle, non documentée & --\\
\vp{mr_k}, \vp{mr_vref} & 2,5 ; $68{,}5\,s$ & -- ; mL & Yoran 1979 (forme) ; calé sur Suzuki 2013 (section~\ref{sec:valves}) & C\\
\vp{k_la}, \vp{k_ra} & 1,45 ; 0,29 & -- & \vp{k_la} : E/A $\approx 1{,}39$, $n = 53$, 10.3390/ani15111524 ; \vp{k_ra} $= 0{,}2\,\vp{k_la}$, placeholder & C / D\\
\vp{ta_atr}, \vp{tra_atr} & 0,09 ; 0,06 & s & non documenté & D\\
\multicolumn{5}{@{}l}{\emph{Autonome et baroréflexe}}\\
\vp{map_set} & 103 & mmHg & PAM 103, Haskins 2005 & A\\
\vp{baro_gain}, \vp{baro_sens} & 1,0 ; 1,0 & -- & non documenté & D\\
\vp{s0}, \vp{v0} & 0,2 ; 0,6 & -- & tonus de référence ($S_\text{REF}$, $V_\text{REF}$ de \cmd{engine.rs}) & D\\
\vp{a_s}, \vp{a_v} & 0,4 ; 0,75 & -- & non documenté & D\\
\vp{rsa_depth} & 0,64 & -- & ASR (pic inspiratoire $-$ creux expiratoire)/FC moyenne $= 40{,}1 \pm 4{,}5$~\%, $n = 23$ chiens vigiles (Grosso 2021, \acc{résumé}) ; modèle 40~\% (0,95 donnait 56~\%). RMSSD 165~ms, dans $259 \pm 120$ (Bogucki 2015) & C\\
\vp{vag_noise}, \vp{sym_noise} & 0,025 ; 0,14 & -- & calage global VFC sur Bogucki 2015 (\acc{lu}) ; \vp{sym_noise} recalé de 0,22 à 0,14 le 18/11/2026 avec le solveur du baroréflexe : LF/HF de la cohorte 0,57 $\to$ 0,31 (cible $0{,}28 \pm 0{,}11$), SDNN, RMSSD et RR dans $\pm 1$~ET & C\\
\multicolumn{5}{@{}l}{\emph{Gaz, acide-base, métabolisme}}\\
\vp{vo2}, \vp{vco2} & $5{,}5\,s^{-0{,}25}$ ; $4{,}6\,s^{-0{,}25}$ & mL/min/kg & VO$_2$ cible 5,4 « dérivée des REE publiées » (PMC6095264, \cmd{targets.json}) ; VCO$_2$ non documentée & C\\
\vp{va0} & 0,0748 & L/min/kg & calé (\cmd{fit\_gas.py}) & C\\
\vp{aagrad}, \vp{edema_k} & 3 ; 4 & mmHg ; mmHg/mmHg & non documenté & D\\
\vp{fio2}, \vp{pb}, \vp{rq} & 0,21 ; 760 ; 0,84 & -- ; mmHg ; -- & constantes usuelles & --\\
\vp{hb} & 15 & g/dL & non documenté & --\\
\vp{p50}, \vp{hill_n} & 30 ; 2,7 & mmHg ; -- & P50 $30{,}0 \pm 1{,}3$, Cambier 2004 (\acc{résumé}) ; $n$ non vérifié & A / --\\
\vp{na}, \vp{k}, \vp{cl} & 150 ; 4,1 ; 113,6 & mmol/L & Constable \& Stämpfli 2005 (cibles 150,7 ; 4,07 ; 113,6 ; \acc{résumé}) & A\\
\vp{alb}, \vp{phos} & 3,77 g/dL ; 1,42 & -- ; mmol/L & idem (37,7 g/L ; 1,42) & A\\
\vp{xa} & $-2{,}48$ & mmol/L & calé (\cmd{fit\_gas.py}) & C\\
\vp{lact0}, \vp{glucose}, \vp{alt} & 1,2 ; 5,2 ; 45 & mmol/L ; mmol/L ; U/L & non documenté & --\\
\multicolumn{5}{@{}l}{\emph{Rein et foie}}\\
\vp{nephron_mass}, \vp{hepatic_fn} & 1 ; 1 & -- & fonctions relatives & --\\
\vp{gfr0} & 3,5 & mL/min/kg & non documenté ; comparé à Bexfield 2008, Kelley 2022, Kelly 2016 (OK, haut de plage) & C\\
\vp{creat0}, \vp{urea0} & 75 ; 5 & \textmu mol/L ; mmol/L & non documenté & --\\
\vp{creat_vd}, \vp{urea_vd} & 0,539 ; 0,66 & L/kg & Collignon 2012 (\acc{résumé}, $n = 50$) ; Bowsher 1985 (\acc{résumé}) & B\\
\multicolumn{5}{@{}l}{\emph{ECG}}\\
\vp{pr}, \vp{qrs}, \vp{qt} & 0,12 ; 0,045 ; 0,20 & s & Lerdweeraphon 2020 (\acc{lu}) ; Bispo 2025 (\acc{lu}) ; Santos 2024 & A\\
\vp{p_amp}, \vp{r_amp}, \vp{t_amp} & 0,2 ; 1,7 ; 0,3 & mV & idem & A\\
\bottomrule
\end{tabularx}

Les paramètres de conduction et de décompensation du même fichier sont donnés aux sections~\ref{sec:rythme} et~\ref{sec:decomp}. Le niveau de preuve par paramètre n'est pas un champ du JSON du patient : la colonne « Év. » reprend ce que disent les notes de provenance et les rapports de vérification, « -- » signifiant qu'aucun document ne l'attribue.

\textbf{Valeurs de sortie du chien sain de 20~kg} (mode déterministe ; \cmd{VERIFICATION\_MALADIES\_MEDICAMENTS.tex}, 18/11/2026, avant l'orifice mitral dynamique) : FC 87,6/min ; PAM 109,4~mmHg ; DC 3,53~L/min ; POG 4,6~mmHg ; PaCO$_2$ 40,2~mmHg ; K$^+$ 4,1~mmol/L ; créatinine 75~\textmu mol/L. Avec $\vp{g_mr} = 0{,}0005$, l'orifice dynamique ne devrait pas changer ces valeurs de façon visible (non recalculé).

\limite{PAM saine recalée à 103 (Haskins) le 18/11/2026, avec la chronologie systolique (elle valait 109,4 depuis $\vp{ts_coef}$ 0,325). Haskins 2005 n'a été lu qu'en résumé (moyennes sans écart-type) ; le DC cible vient de l'IC par une formule de surface corporelle. Population de référence hétérogène (races, méthodes). Calage, et non validation prospective : aucune donnée n'est tenue à part.}

%=====================================================================
\section{Cœur : ventricules gauche et droit}\label{sec:coeur}
%=====================================================================

\textbf{Rôle.} Pompe à élastance variable dans le temps, une pour chaque ventricule ; relation télésystolique linéaire, relation télédiastolique exponentielle.

\subsection{Activation et durée de systole}
Rythme régulier (\cmd{hemo::activation\_g}) : période $T = 60/\text{FC}$, $t_c = t \bmod T$, systole $t_s = \vp{ts_coef}\sqrt{T}$, relaxation $t_r = \vp{tr_frac}\,t_s$,
\begin{equation}
e(t) = \begin{cases}
\tfrac12\bigl(1 - \cos(\pi (t_c/t_s)^{\gamma})\bigr) & 0 \le t_c < t_s,\\
\tfrac12\bigl(1 + \cos(\pi (t_c - t_s)/t_r)\bigr) & t_s \le t_c < t_s + t_r,\\
0 & \text{sinon.}
\end{cases}
\end{equation}
avec $\gamma = \vp{act_gamma}$ : $\gamma < 1$ raidit le début de la contraction. L'origine $t_c = 0$ est le \emph{début mécanique} : début du QRS tracé par l'ECG + \vp{emd_v} (\cmd{hemo::mech\_offset}) ; la contraction atriale débute \vp{emd_a} après le début de l'onde P.

\textbf{Recalage de la chronologie (18/11/2026).} Le cœur animé a montré une contraction isovolumique de 90~ms, une période pré-éjectionnelle (PEP, du début du QRS à l'ouverture aortique) de 110 à 136~ms, une relaxation isovolumique (IVRT) de 88~ms et un indice de Tei de 1,0, avec un d$P$/d$t_\text{max}$ de 1\,670~mmHg/s. Avec $\vp{ts_coef} = 0{,}25$ (0,325 auparavant), $\gamma = 0{,}55$ et $\vp{tr_frac} = 0{,}35$ (0,4), puis le recalage global de \vp{r_sys}, \vp{vol_shift}, \vp{ees_lv}, \vp{r_pul} et \vp{b_lv} sur les mêmes cibles que \cmd{fit\_hemo.py} (PAM 103, IC 4,42~L/min/m$^2$, PAP 14, POG 5,5, FE 0,55), le chien sain de 20~kg à 87/min donne les valeurs du tableau~\ref{tab:sti}. PAM 103,0 (109,4 auparavant), d$P$/d$t_\text{max}$ 2\,840~mmHg/s.

\begin{table}[h]\centering\small
\caption{Intervalles du cycle cardiaque gauche et droit, chien sain : modèle et références.}\label{tab:sti}
\begin{tabular}{@{}lll@{}}
\toprule
Intervalle & Modèle & Référence (chien, accès)\\
\midrule
PEP aortique & 54 ms & $58{,}8 \pm 4{,}8$ (SE), $n = 12$, vigiles (Trikhun 2020, \acc{lu}) ; $33 \pm 26$, $n = 14$ (Gicana 2022, \acc{lu})\\
Temps d'éjection VG & 178 ms & $170{,}6 \pm 18{,}4$ à 102/min (Gicana 2022) ; 173 (168--179), $n = 45$ (Baumwart 2005, \acc{résumé})\\
IVRT & 61 ms & $61{,}4 \pm 9{,}7$, $n = 49$, vigiles (Favier 2025, \acc{lu})\\
Contraction isovolumique & 35 ms & --\\
Indice de Tei VG & 0,54 & $0{,}42 \pm 0{,}22$ (Gicana) ; 0,41 (Baumwart) ; $0{,}54 \pm 0{,}20$ (Guglielmini 2019, \acc{lu})\\
Fermeture $\to$ ouverture mitrale & 273 ms & $239{,}7 \pm 34{,}2$ à 102/min (Gicana)\\
PEP et éjection pulmonaires & 33 ; 204 ms & $52{,}8 \pm 2{,}0$ (SE) ; $147 \pm 8$ (Trikhun) ; éjection VD 187 (182--192) (Baumwart)\\
\bottomrule
\end{tabular}
\end{table}

\subsection{Pression ventriculaire}
\begin{equation}
P_v = e_\text{act}\,E_{es}\,(V - V_0) + (1 - e)\,a\left(\exp\bigl(\mathrm{clip}(b\,(V - V_0), -20, 8)\bigr) - 1\right),
\end{equation}
avec $(E_{es}, V_0, a, b) = (\vp{ees_lv}, \vp{v0_lv}, \vp{a_lv}, \vp{b_lv})$ pour le VG et les homologues \vp{_rv} pour le VD. En rythme régulier $e_\text{act} = e$.

\subsection{Rythme irrégulier : battement par battement}
Quand la conduction n'est pas sinusale 1:1 (section~\ref{sec:rythme}), chaque battement $i$ a sa systole et son amplitude (\cmd{hemo::beat\_schedule}) :
\begin{align}
t_{s,i} &= \vp{ts_coef}\sqrt{T_0}\,\bigl(RR_i/T_0\bigr)^{\vp{ts_restitution}}, \qquad t_{r,i} = \vp{tr_frac}\,t_{s,i},\\
x &= \exp\bigl(-\max(0, RR_i - t_{s,i-1})/\vp{ca_tau}\bigr),\quad \text{rel} = u\,(1 - x),\quad u \leftarrow u\,x + \rho\,\text{rel} + (1 - \rho),\\
c_i &= \text{rel}\times\begin{cases}\vp{ees_wide} & \text{QRS large}\\ 1 & \text{sinon}\end{cases},\qquad u_0 = \frac{1}{1 - \exp(-(T_0 - t_{s,0})/\vp{ca_tau})},
\end{align}
avec $\rho = \vp{ca_recirc}$ et $T_0$ la période moyenne. C'est un modèle de recirculation du Ca$^{2+}$ à deux compartiments : en régime établi $c = 1$ à toute fréquence ; un battement prématuré libère moins (restitution), le suivant plus (potentialisation post-extrasystolique). À l'instant $t$, $e$ est le maximum des activations des deux derniers battements et $e_\text{act}$ le maximum des $c_j\,e_j$.

\sources{Structure du modèle calcique : Schouten 1987, Biophys J, 10.1016/s0006-3495(87)83307-6 ; principe : Wier 1986, J Physiol, 10.1113/jphysiol.1986.sp016167 ; potentialisation : Cooper 1993, Circulation, 10.1161/01.cir.88.6.2962 ; dyssynchronie du QRS large : Prinzen 1999, JACC, 10.1016/s0735-1097(99)00068-6 (DOI vérifiés dans OpenAlex, accès non précisé). \vp{ts_coef} : PMC9394128 (beagle, $n = 14$). \vp{ees_lv}, \vp{b_lv} : calés globalement (\cmd{fit\_hemo.py}).}

\textbf{Statut.} \vp{ca_recirc}, \vp{ts_restitution}, \vp{emd_a} : \evD{}. \vp{ts_coef}, \vp{act_gamma}, \vp{tr_frac}, \vp{ca_tau}, \vp{ees_wide} : \evC{}. \vp{emd_v} : \evB{}.

\textbf{Extrasystoles (vérification du 18/11/2026).} Sur 19 ESV isolées (patient à ESV, sévérité 0,8, $6 \times 30$~s) :
\begin{itemize}
\item \textbf{ECG.} QRS de 82~ms (bornes $\pm 2\sigma$), contre 44~ms en sinusal ; onde T discordante (environ $-0{,}5$~mV, contre $+0{,}3$).
\item \textbf{Volume d'éjection selon le couplage.} Nul sous environ 300~ms de couplage, puis 21 à 32~\% du VES sinusal vers 380--410~ms et 45 à 70~\% de 450 à 640~ms. Le déficit de pouls est bien présent.
\item \textbf{Potentialisation.} Le battement suivant est potentialisé : $c$ jusqu'à 1,34, VES jusqu'à $+36$~\%.
\end{itemize}
Repère : le VES de l'ESV est de 2,5 à 4,1~mL à 200~ms de couplage, et de 10,2 à 14,8~mL à 375~ms, pour 16,4 à 17,3~mL en stimulation régulière à 400~ms (Potfay 2015, PMC4618025, \acc{lu}). Ces chiens étaient anesthésiés, thorax ouvert : la pression aortique y est plus basse que chez le chien vigile modélisé, ce qui favorise l'éjection précoce. Le modèle reste en dessous de ces valeurs entre 300 et 380~ms.

\limite{Le modèle à élastance ne représente ni la précharge recrutable de Frank--Starling au-delà de la relation $E_{es}$, ni la relaxation active (lusitropie) : le sympathique et le pimobendane n'agissent que sur $E_{es}$. La pause compensatrice après une ESV est souvent incomplète et les ESV interpolées sont fréquentes, car le nœud AV ne reste réfractaire que \vp{av_erp} (0,15~s, calé sur la FA) et la pénétration rétrograde de l'ESV dans la jonction n'est pas représentée. Aucune donnée canine chiffrée n'a été trouvée sur la fréquence des pauses complètes ; selon Shoureshi 2023 (PMC10513922, \acc{lu}), le chien ne semblerait pas conduire de façon rétrograde vers les oreillettes pendant les ESV. Le gradient VD--AP en systole atteint 8~mmHg (résistance valvulaire commune aux quatre valves).}

%=====================================================================
\section{Circulations systémique et pulmonaire, volémie}\label{sec:circ}
%=====================================================================

\textbf{Rôle.} Six compartiments en série, états = volumes $V = [V_{lv}, V_{rv}, V_{ao}, V_{sv}, V_{pa}, V_{pv}]$ : VG, VD, aorte, veines systémiques (avec l'OD), artère pulmonaire, veines pulmonaires (avec l'OG).

\subsection{Pressions vasculaires}
\begin{align}
P_{ao} &= \frac{V_{ao} - \vp{v0_ao}}{\vp{c_ao}}, &
P_{sv} &= (1 + \vp{k_ra}\,e_a)\,\frac{V_{sv} - \vp{v0_sv}}{\vp{c_sv}\cdot\vp{ven_cap}},\\
P_{pa} &= \frac{V_{pa} - \vp{v0_pa}}{\vp{c_pa}}, &
P_{pv} &= (1 + \vp{k_la}\,e_a)\,\frac{V_{pv} - \vp{v0_pv}}{\vp{c_pv}}.
\end{align}
$e_a$ est l'activation atriale (section~\ref{sec:oreillettes}). \vp{ven_cap} (capacitance veineuse relative, 1 = basale) n'intervient que dans la pression, pas dans le volume initial : une veinodilatation redistribue le volume sans en ajouter. Il a été ajouté le 18/11/2026 pour la veinodilatation du furosémide.

\subsection{Débits et bilan}
\begin{align}
Q_{sys} &= \frac{P_{ao} - P_{sv}}{\vp{r_sys}}, \qquad Q_{pul} = \frac{P_{pa} - P_{pv}}{\vp{r_pul}},\\
\dot V_{lv} &= Q_{mi} - Q_{ao}, & \dot V_{rv} &= Q_{tr} - Q_{pu}, & \dot V_{ao} &= Q_{ao} - Q_{sys},\\
\dot V_{sv} &= Q_{sys} - Q_{tr}, & \dot V_{pa} &= Q_{pu} - Q_{pul}, & \dot V_{pv} &= Q_{pul} - Q_{mi},
\end{align}
les débits valvulaires $Q_{mi}, Q_{ao}, Q_{tr}, Q_{pu}$ étant donnés à la section~\ref{sec:valves}. Le volume total est conservé.

\subsection{Volumes initiaux et volémie}
\cmd{initial\_volumes} : ventricules sur leur relation télédiastolique à 8 (VG) et 4~mmHg (VD), vaisseaux à des pressions de départ usuelles, et réservoir veineux systémique porteur de la volémie :
\begin{align}
V_{lv}^0 &= \vp{v0_lv} + \ln(1 + 8/\vp{a_lv})/\vp{b_lv}, \qquad V_{rv}^0 = \vp{v0_rv} + \ln(1 + 4/\vp{a_rv})/\vp{b_rv},\\
V_{ao}^0 &= \vp{v0_ao} + 100\,\vp{c_ao}, \quad V_{pa}^0 = \vp{v0_pa} + 14\,\vp{c_pa}, \quad V_{pv}^0 = \vp{v0_pv} + 6\,\vp{c_pv},\\
V_{sv}^0 &= \vp{v0_sv} + 4\,\vp{c_sv} + (\vp{bvf} - 1)\cdot 85\,\vp{bw} + \vp{vol_shift}.
\end{align}
La volémie de référence est donc 85~mL/kg (source non documentée) ; \vp{bvf} la multiplie, \vp{vol_shift} (mL) l'augmente ou la diminue. Le baroréflexe et le diurétique agissent sur \vp{vol_shift}.

\subsection{Simulation et indicateurs}
Régime périodique : 11 battements (\vp{N_BEATS}), indicateurs sur le dernier, échantillonné à $\Delta t = 2$~ms. PAM, PAS, PAD = moyenne, maximum, minimum de $P_{ao}$ ; DC $= \overline{Q_{sys}}\times 60/1000$~L/min ; VTD, VTS = max et min de $V_{lv}$ ; FE $= (\text{VTD} - \text{VTS})/\text{VTD}$ ; POG moyenne et maximale = $P_{pv}$ ; PAP moyenne ; PVC $= \overline{P_{sv}}$ ; PTDVG = $P_{lv}$ au début du battement ; \vp{converged} si le volume d'éjection varie de moins de 3~\% entre les deux derniers battements. Rythme irrégulier : 30~s de simulation, moyennes sur les battements complets après 10~s (\cmd{simulate\_hemo\_beats}).

\sources{Paramètres \vp{r_sys}, \vp{r_pul}, \vp{vol_shift} calés sur Haskins 2005 (PAM 103, PAP 14, POG 5,5, IC 4,42 ; \acc{résumé}). Compliances et volumes non contraints : non documentés. Volémie 85~mL/kg : non documentée.}

\textbf{Statut.} \evA{} pour les cibles du sain (moyennes publiées) ; \evD{} ou non attribué pour les compliances.

\limite{Compliances linéaires et résistances constantes (pas d'autorégulation locale, pas de répartition régionale des débits ni de perfusion coronaire explicite : la pression de perfusion coronaire n'entre que par le critère $\text{PAD} - \text{PTDVG}$ de la section~\ref{sec:decomp}). Pas de compartiment interstitiel : le remplissage capillaire après diurèse est résumé par $\tau_\text{refill}$ (section~\ref{sec:furo}). En MMVD, faire baisser la POG demande un grand déficit de volume, car la POG dépend surtout de \vp{c_pv} et de la régurgitation (correction 7 du 18/11, d'abord « non appliquée », puis traitée par l'orifice dynamique et \vp{ven_cap}, section~\ref{sec:valves}).}

%=====================================================================
\section{Valves}\label{sec:valves}
%=====================================================================

\textbf{Rôle.} Quatre diodes lisses avec conductance de fuite rétrograde ; la fuite mitrale porte la MMVD.

\subsection{Diode lisse}
\begin{equation}
\sigma = \tfrac12\bigl(1 + \tanh(\Delta p/0{,}3)\bigr), \qquad Q = \sigma\,\frac{\Delta p}{\vp{r_valve}} + (1 - \sigma)\,g\,\Delta p,
\end{equation}
avec $\Delta p = P_{pv} - P_{lv}$ (mitrale), $P_{lv} - P_{ao}$ (aortique), $P_{sv} - P_{rv}$ (tricuspide), $P_{rv} - P_{pa}$ (pulmonaire), et $g = g_\text{eff}$ (mitrale), \vp{g_ar}, \vp{g_tr}, \vp{g_pr}, multipliées par $(\vp{bw}/20)^{0{,}75}$ : les conductances de fuite sont normalisées au chien de référence de 20~kg et suivent les débits, comme $1/\vp{r_valve}$ (\cmd{hemo::leak\_scale}, 18/11/2026 ; auparavant absolues, si bien qu'au milieu du stade C la POG valait 37~mmHg à 3~kg et 8 à 60~kg). La largeur 0,3~mmHg du \emph{tanh} est fixée dans le code.

\subsection{Orifice mitral dynamique (18/11/2026)}
\begin{equation}
g_\text{eff} = \vp{g_mr}\left(\frac{\text{VTD}}{\vp{mr_vref}}\right)^{\vp{mr_k}}, \qquad \vp{mr_k} = 2{,}5,\quad \vp{mr_vref} = 68{,}5\,s~\text{mL}
\end{equation}
L'orifice est \emph{fixe pendant le battement} et suit la taille du VG d'un battement à l'autre : \cmd{engine::hemo\_with\_rhythm} simule l'hémodynamique avec un orifice statique, lit le VTD, puis itère un point fixe amorti sur $\ln g$ (relaxation $\tfrac12$, tolérance $10^{-3}$, au plus 12 simulations). Il est ignoré si $\vp{g_mr} < 0{,}05$ (chien sain). La conductance retenue est reportée dans les paramètres de l'évaluation (\vp{g_mr}, $\vp{mr_k} = 0$) pour que l'ECG et le son recalculés en aval restent cohérents. Une première version (18/11) faisait dépendre l'orifice du volume \emph{instantané} : le pic de débit régurgité en début de systole était multiplié par 4,7 (souffle de grade 6 à MMVD 0,3) ; la forme par battement le ramène à 46~mL/s (grade 3, contre 33~mL/s avant l'orifice dynamique). Une baisse de précharge (diurétique, veinodilatation) réduit l'orifice et la fraction régurgitée, ce qui permet à la POG de baisser. \vp{mr_vref} correspond au VTD du chien sain (68,5~mL à 20~kg depuis le recalage du 18/11/2026, 72 auparavant).

\sources{Forme : Yoran C, Yellin EL, Becker RM et al. 1979, PMID 445720 (6 chiens à thorax ouvert, \acc{résumé} : l'expansion volémique augmente surtout l'orifice régurgitant effectif ; la noradrénaline, qui réduit le VG, le diminue). Exposant : calé sur Suzuki 2013 (IM expérimentale, furosémide 0,17~mg/kg/h pendant 6~h : POG 14,75 $\to$ 10,77~mmHg, $-27$~\% ; \acc{résumé}, J Vet Intern Med, doi:10.1111/jvim.12154 non revérifié) ; le modèle donnait 13,66 $\to$ 9,95~mmHg ($-27$~\%) ; après le recalage de la congestion, 14,92 $\to$ 11,21~mmHg ($-25$~\%) à MMVD 0,72, sévérité dont la POG de base égale celle de l'étude. Le test \cmd{furosemide\_lowers\_lap\_in\_mitral\_regurgitation} (\cmd{tests/drug\_corrections.rs}) exige une baisse d'au moins 18~\% de la POG à MMVD 0,78 avec une PAM conservée à 90~\%, et vérifie qu'avec orifice fixe et sans veinodilatation la POG baisse de moins de 10~\%.}

\subsection{Fraction de régurgitation}
\begin{equation}
\text{RF} = \frac{\Delta t\left(-\sum \min(Q_{mi}, 0) - \sum \min(Q_{ao}, 0)\right)}{\text{VTD} - \text{VTS}},
\end{equation}
sur le dernier battement (ou la somme des battements en rythme irrégulier).

\textbf{Statut.} Fuite de base (\vp{g_mr} = 0,0005 etc.) : non documentée. Orifice dynamique : \evC{} (forme sourcée, exposant calé sur un résumé).

\limite{Une seule fuite par valve, sans sténose (\vp{r_valve} identique pour les quatre valves). Les coefficients de la MMVD sur \vp{g_mr} ont été recalés le 18/11/2026 pour tenir compte de l'orifice dynamique (section~\ref{sec:mmvd}) ; la fraction de régurgitation par stade n'a pas été recomparée à Kawai 2026 depuis.}

%=====================================================================
\section{Oreillettes}\label{sec:oreillettes}
%=====================================================================

\textbf{Rôle.} Contribution atriale au remplissage (onde A). L'OG est incluse dans le compartiment \vp{pv}, l'OD dans \vp{sv} ; leur pression passive est multipliée pendant l'activation atriale (équations de la section~\ref{sec:circ}).

\textbf{Activation.} Rythme régulier : $e_a(t) = \mathrm{act}(t + \text{PR}, T, \vp{ta_atr}, \vp{tra_atr})$, la même fonction demi-cosinus que $e(t)$ avec $t_s = \vp{ta_atr} = 0{,}09$~s et $t_r = \vp{tra_atr} = 0{,}06$~s : l'oreillette se contracte PR secondes avant le ventricule. Rythme irrégulier : activation à chaque activation atriale \emph{organisée} de la conduction, aucune en fibrillation ou en arrêt atrial. Avec $\vp{k_la} = \vp{k_ra} = 0$, le modèle redevient le modèle passif d'origine au bit près.

\textbf{Sorties.} \vp{mitral_e} et \vp{mitral_a} : pics du débit mitral positif hors et pendant l'activation atriale.

\sources{$\vp{k_la} = 1{,}45$, calé sur E/A $\approx 1{,}39$ (E 0,82, A 0,59~m/s, $n = 53$ ; 10.3390/ani15111524, accès non précisé) à la FC de repos du modèle ; contribution atriale $\approx 10$~\% du volume d'éjection. $\vp{k_ra} = 0{,}2\,\vp{k_la}$ : placeholder D. \vp{ta_atr}, \vp{tra_atr} : non documentés.}

\limite{Les ondes E et A fusionnent au-delà d'environ 120/min. La pompe atriale a abaissé la POG en MMVD (à $s = 0{,}8$ : 20,6 $\to$ 17,4~mmHg) : l'œdème n'apparaît plus qu'au-delà de $s \approx 0{,}85$. Pas de volume atrial séparé, donc pas de dilatation atriale mesurable (LA:Ao).}

%=====================================================================
\section{Système autonome et baroréflexe}\label{sec:auto}
%=====================================================================

Deux niveaux : un baroréflexe \emph{stationnaire} dans l'évaluation (\cmd{engine::evaluate\_once}), et une simulation \emph{dynamique} à 20~Hz qui produit la VFC (\cmd{autonomic.rs}).

\subsection{Baroréflexe stationnaire}
Commande barorécepteur $d$ (positive en hypotension), initialisée à 0. À chaque passage (au plus 5) :
\begin{align}
s_\text{tot} &= \mathrm{clip}(\vp{s0} + 0{,}8\,d,\;0, 1), \qquad v_\text{tot} = \mathrm{clip}(\vp{v0} + \vp{vagal_add} - 0{,}6\,d,\;0, 1),\\
s_\text{eff} &= s_\text{tot}\,(1 - \vp{beta_block}), \qquad v_\text{eff} = v_\text{tot}\,(1 - \vp{vagal_block}),\\
\text{FC}_\text{sinus} &= \vp{hr}\cdot\frac{h(s_\text{eff}, v_\text{eff})}{h(0{,}2;\,0{,}6)}, \qquad h(s, v) = (1 + \vp{a_s}\,s)(1 - \vp{a_v}\,v),\\
E_{es}' &= E_{es}\,\bigl(1 + 0{,}4\,(s_\text{eff} - \vp{s0})\bigr)\quad\text{(VG et VD)},\\
R_{sys}' &= \vp{r_sys}\,\bigl(1 + 0{,}5\,(s_\text{tot} - \vp{s0})\bigr), \qquad
\vp{vol_shift}' = \vp{vol_shift} + 0{,}20\cdot 85\,\vp{bw}\,(s_\text{tot} - \vp{s0}),\\
d_t &= \mathrm{clip}\bigl(\vp{baro_gain}\,(\vp{map_set} - \text{PAM})/45,\;-1, 1\bigr).
\end{align}
\textbf{Solveur du point fixe} (18/11/2026, \cmd{engine::baro\_next}). On résout $g(d) = d_t(d) - d = 0$ sur $[-1, 1]$ : $g$ décroît avec $d$ et $g(-1) \ge 0 \ge g(1)$, la racine est unique et encadrée. Le premier pas est amorti ($d + g/2$, l'ancienne itération), puis la sécante est appliquée sur les deux derniers points, avec bisection de l'encadrement courant si elle en sort. Arrêt si $|g| < 0{,}01$ (environ 0,45~mmHg), au plus 16 évaluations. L'itération amortie précédente (arrêt à $|g| < 0{,}04$, au plus 5 passages) s'arrêtait dans certains cas sur un état non convergé et incohérent. Sous opioïde, elle donnait une PAM de 88~mmHg avec un sympathique en baisse ; le point fixe donne 102~mmHg, FC $-32$~\%, conforme à Mills 2025 et à Maiante 2009. En MMVD C, elle donnait PAM 124 et FC 114, contre 114 et 105 au point fixe. Le traitement sur 7 jours est aussi plus rapide (4,6~s contre 6,8~s en natif). Le bêtabloquant n'agit que sur la FC, l'inotropisme et la conduction AV ($s_\text{eff}$), pas sur la vasoconstriction ni le recrutement veineux ($s_\text{tot}$). La FC rapportée est la FC \emph{ventriculaire} ($60/T$ de l'hémodynamique), la FC sinusale reste disponible (\vp{sinus_hr}).

\textbf{Effet dromotrope du sympathique} (18/11/2026), au même passage :
\begin{equation}
\phi = \max\bigl(0{,}1;\;1 - \vp{av_symp}\,(s_\text{eff} - 0{,}2)\bigr),\qquad
\vp{av_erp}' = \vp{av_erp}\,\phi,\qquad \text{PR}' = \vp{hv} + (\vp{pr} - \vp{hv})\,\phi,
\end{equation}
avec $\vp{av_symp} = 0{,}5$ : le sympathique raccourcit la période réfractaire et le délai AH du nœud AV ; le bêtabloquant, qui réduit $s_\text{eff}$, ralentit donc la réponse ventriculaire de la FA.

\textbf{Modulation autonome des foyers d'échappement} (18/11/2026), au même passage :
\begin{equation}
\vp{junct_rate}' = \vp{junct_rate}\,\frac{1 + \vp{esc_a_s}\,s_\text{eff}}{1 + 0{,}2\,\vp{esc_a_s}}\cdot\frac{\max(0,\,1 - \vp{esc_a_v}\,v_\text{eff})}{1 - 0{,}6\,\vp{esc_a_v}},\qquad
\vp{vent_escape_rate}' = \vp{vent_escape_rate}\,\frac{1 + \vp{esc_a_s}\,s_\text{eff}}{1 + 0{,}2\,\vp{esc_a_s}}.
\end{equation}
Le foyer jonctionnel est accéléré par le sympathique et ralenti par le vague ; le foyer ventriculaire n'est sensible qu'au sympathique. Chez le chien vigile, l'atropine raccourcit la pause d'échappement jonctionnel (4,5 $\to$ 3,4~s, soit $\times 1{,}32$ sur la fréquence ; réflexes intacts, d'où un calage en boucle fermée avec le baroréflexe convergé : $\vp{esc_a_v} = 0{,}72$, qui donne $\times 1{,}32$ en bloc sino-atrial complet) et le practolol l'allonge (4,9 $\to$ 9,4~s ; Tanahashi 1979, \acc{résumé}). Cette pause mêle fréquence et suppression par entraînement : $\vp{esc_a_s} = 2$ n'en attribue qu'une partie à la fréquence ($\times 0{,}71$ sous blocage complet). Un gain plus fort (4,6, tout attribué à la fréquence) rend le point fixe du baroréflexe instable en bloc sino-atrial complet. \evC{}

\textbf{Suppression de \vp{kfac}.} Le frein multiplicatif de K$^+$ sur la FC ($\max(0{,}4; 1 - 0{,}30\max(0, K - 6))$) a été supprimé le 17/11/2026 : K$^+$ agit désormais sur la conduction (section~\ref{sec:rythme}).

\sources{\vp{map_set} = 103 (Haskins 2005). Gains (0,8 ; 0,6 ; 0,4 ; 0,5 ; 0,20 ; 45), \vp{a_s}, \vp{a_v}, \vp{s0}, \vp{v0} : non documentés (placeholders). \vp{av_symp} : calé sur la dénervation sympathique en FA, FC ventriculaire 144,5 $\to$ 121,5/min ($-16$~\%) (PMC7872775, \acc{lu} ; Rigel 1990, \acc{résumé}) ; \evC{}.}

\subsection{Simulation autonome dynamique (VFC)}
Pas $\Delta t = 1/20$~s, 20~s de mise en régime puis 120~s enregistrées (par défaut). Entrées : FC sinusale, $\bar s = s_\text{eff}$, $\bar v = v_\text{eff}$, PAM $\bar m$, fréquence respiratoire $f_r$ et ventilation alvéolaire VA de l'évaluation.
\begin{itemize}
\item Respiration : $\dd\xi = -\xi\,\Delta t/8 + 0{,}05\sqrt{\Delta t}\,\mathcal N$ ; phase $\varphi_r \mathrel{+}= f_r(1 + \xi)\,\Delta t$ ; volume pulmonaire $\rho(\varphi)$ = demi-cosinus montant sur 40~\% du cycle (inspiration), descendant ensuite.
\item Profondeur ventilatoire relative $\nu = (\text{VA}/\text{FR})/(\vp{va0}\,\vp{bw}/20)$ ; variation respiratoire de PA $a_{ip} = 2\nu$~mmHg ; profondeur de RSA $m = \vp{rsa_depth}\cdot\mathrm{clip}(\nu^{0{,}7}, 0{,}3, 1{,}5)$.
\item PAM instantanée : $m_p = \bar m - a_{ip}\bigl(\rho(\varphi_r - 0{,}25) - 0{,}5\bigr) + 25\,(s_e - \bar s)$.
\item Barorécepteur sigmoïde $f_b(m) = 1/(1 + e^{-(m - \vp{map_set})/10})$ ; retards 2,5~s (sympathique) et 0,3~s (vagal) ; $\Delta f_{s,v} = f_b(m_p(t - \text{retard})) - f_b(\bar m)$.
\item Cibles : $s^* = \bar s - 1{,}4\,\vp{baro_sens}\,\Delta f_s$ ; $v^* = \bar v\,(1 - m(\rho - 0{,}5)) + \vp{baro_sens}\,\Delta f_v + x_v$, avec la dérive vagale $\dd x_v = -x_v\Delta t/12 + \vp{vag_noise}\sqrt{\Delta t}\,\mathcal N$.
\item Dynamiques : $s_e \mathrel{+}= \frac{\Delta t}{3}(s^* - s_e) + \vp{sym_noise}\sqrt{\Delta t}\,\mathcal N$ ; $v_e \mathrel{+}= \frac{\Delta t}{0{,}25}(v^* - v_e)$ ; bornes $s \in [0; 1{,}2]$, $v \in [0; 1]$.
\item FC instantanée $= \mathrm{clip}\bigl(\overline{\text{FC}}\,h(s, v)/h(\bar s, \bar v), 20, 300\bigr)$ ; nœud sinusal à intégration de phase (IPFM) : une onde R quand la phase cumulée $\int \text{FC}/60\,\dd t$ atteint 1 (instant interpolé).
\end{itemize}
\textbf{Indices de VFC} (\cmd{hrv\_metrics}, au moins 20 RR) : RR moyen, SDNN, RMSSD ; RR rééchantillonné à 4~Hz, densité spectrale de Welch (fenêtre de Hann, segments de $\min(256, n)$ points) ; LF 0,04--0,15~Hz, HF 0,15--1~Hz (bandes adaptées au chien) ; amplitude de RSA en battements/min par cycle respiratoire.

\sources{Calage global (non paramètre par paramètre) sur Bogucki \& Noszczyk-Nowak 2015, Pol J Vet Sci 18:307, doi:10.1515/pjvs-2015-0040 (50 chiens sains, Holter 60~min ; \acc{lu}) : RR moyen $677{,}7 \pm 126{,}9$~ms, SDNN $208{,}9 \pm 77{,}1$, RMSSD $259{,}0 \pm 120{,}2$, LF/HF $0{,}28 \pm 0{,}11$ (\cmd{calibration/targets.json}, validation statistique sur $\ge 20$ graines, tolérance $\pm 1$~ET). Sens de la relation VFC--FC : Billman 2013, Front Physiol 4:222, doi:10.3389/fphys.2013.00222 (\acc{lu}). Constantes de temps, retards, largeur de sigmoïde, gain vasomoteur 25~mmHg : non documentés (\evD{}).}

\limite{Stochastique : validation statistique, pas au point. LF/HF de la cohorte saine 0,31 (cible $0{,}28 \pm 0{,}11$) depuis le recalage du 18/11/2026 : l'écart déclaré du 17/11 est levé, les quatre indicateurs tiennent $\pm 1$~ET dans \cmd{tests/signal\_properties.rs}. Source VFC unique, à valeurs élevées et partiellement incohérente (FC de l'article 102 contre NN moyen 677~ms). En MMVD C : SDNN trop bas ($-54$~ms), LF/HF sans hausse (Bogucki 2017). RSA trop ample en tachycardie (modulation relative du tonus vagal) ; RSA nulle si la fréquence respiratoire dépasse $\approx 70$/min (échantillonnage par battements). Pas de halètement. La simulation dynamique n'agit pas en retour sur l'hémodynamique stationnaire.}

%=====================================================================
\section{Conduction et rythme}\label{sec:rythme}
%=====================================================================

\textbf{Rôle.} Générer et propager l'influx battement par battement (\cmd{rhythm.rs}) : nœud sinusal, oreillettes, nœud AV, His--Purkinje, échappements, foyer ectopique, fibrillation atriale, effets de K$^+$. FC, ECG, oreillettes et sons en découlent ; les modules de troubles du rythme n'agissent que sur ces paramètres.

\subsection{Nœud AV}
\begin{align}
\nu_v &= 1 + \vp{av_vagal}\,(v_\text{eff} - 0{,}6), \qquad \kappa = 1 + \vp{av_k_gain}\,\max(0, K - 5{,}5),\\
\text{AH}_0 &= \max(0, \vp{pr} - \vp{hv})\,\nu_v\,\kappa, \qquad \text{ERP} = \vp{av_erp}\,\nu_v\,\kappa,\\
\text{AH} &= \text{AH}_0 + \vp{av_beta}\,\exp\bigl(-(r - \text{ERP})/\vp{av_tau}\bigr),
\end{align}
où $r = t - t_\text{His}$ est le temps de récupération depuis la dernière traversée nodale ($\vp{pr}$ et $\vp{av_erp}$ sont ceux déjà modulés par $\phi$, section~\ref{sec:auto}). Si $r < \text{ERP}$ (ou si un influx est déjà en cours), l'influx est bloqué ; s'il arrive en réfractarité relative ($r \ge 0{,}5\,\text{ERP}$), il est \emph{caché} : $t_\text{His} \mathrel{+}= \vp{av_conceal}\cdot r$. Sinon $t_\text{His} = t + \text{AH}$ et l'activation ventriculaire est prévue à $t + \text{AH} + \vp{hv}$ (PR du battement $= \text{AH} + \vp{hv}$).

\subsection{His--Purkinje, ventricule, échappements, ectopie}
\begin{itemize}
\item Conduction infranodale avec probabilité $p_{HP} = \vp{hp_conduction}\times\bigl(1$ si $K < \vp{k_hp_block}$, sinon $\max(0, 1 - \vp{hp_k_slope}(K - \vp{k_hp_block}))\bigr)$.
\item Période réfractaire ventriculaire $= \vp{qt}$ : un influx arrivant plus tôt après la dernière activation est perdu.
\item Échappements jonctionnel (QRS fin) et ventriculaire (QRS large) : cycle $60/(\text{taux}\cdot k_\text{esc})$, $k_\text{esc} = \max(0, 1 - \vp{esc_k_gain}\max(0, K - 5{,}5))$, compté depuis la dernière activation ventriculaire (fréquences modulées par le système autonome, section~\ref{sec:auto}). L'échappement jonctionnel exige $p_{HP} > 0$.
\item Suppression par entraînement (overdrive suppression) : le cycle d'un foyer est multiplié par $1 + \vp{esc_sup_gain}\max(0,\ \text{cycle}\cdot\bar f - 1)$, où $\bar f$ est la fréquence ventriculaire moyennée sur \vp{esc_sup_tau} : à chaque activation (toute origine) d'intervalle $\Delta t$, $\bar f \mathrel{+}= (1 - e^{-\Delta t/\vp{esc_sup_tau}})(1/\Delta t - \bar f)$. Un patient qui conduit part de $\bar f = 1/T_\text{sinus}$ (entraîné de longue date) ; en bloc complet ou sans influx sinusal, de $\bar f = 0$. La suppression croît avec la fréquence et la durée de l'entraînement (charge de la pompe Na/K) et s'efface en quelques dizaines de battements d'échappement. Chien sain à 90/min : premier échappement jonctionnel 3,3~s après l'arrêt du sinus (2,6 cycles), puis cycles qui se raccourcissent vers 1,28~s.
\item Foyer ectopique ventriculaire : processus de Poisson de taux $\vp{vpc_rate} + \vp{vpc_k_gain}\max(0, \vp{k_vpc} - K)$ /min (accru en hypokaliémie) ; une ESV rend le nœud AV réfractaire (conduction rétrograde).
\end{itemize}

\subsection{Oreillettes et nœud sinusal}
\begin{itemize}
\item Rythme sinusal : décharges à la période sinusale (ou aux instants du modèle autonome pour les signaux) ; chaque décharge est bloquée en sortie avec probabilité \vp{sa_block}.
\item Fibrillation ($\vp{af} \ge 0{,}5$) : activations désorganisées d'intervalle $\max\bigl(0{,}02;\ \vp{af_cycle}(1 + \vp{af_cv}\sqrt3(2u - 1))\bigr)$, $u$ uniforme (coefficient de variation \vp{af_cv}).
\item Arrêt atrial ($K \ge \vp{k_atrial_standstill}$) : plus d'activation atriale (ni onde P ni contraction), mais conduction sino-ventriculaire conservée.
\end{itemize}

\subsection{Couplage avec l'hémodynamique}
Si le rythme est sinusal 1:1 à PR stable (40 battements simulés sans aléa, 10 derniers PR identiques à $10^{-6}$ près ; aucun processus aléatoire ni arrêt atrial), l'hémodynamique reste périodique avec le PR de la conduction. Sinon : séquence de conduction sur 30~s (aléa SplitMix64, graine 0, identique en Python et en Rust), puis hémodynamique battement par battement, moyennes après 10~s. Moins de 3 battements après 10~s : décès \vp{cardiac_arrest}.

\begin{tabularx}{\linewidth}{@{}P{3.3cm}P{1.6cm}LP{0.7cm}@{}}
\caption{Paramètres de conduction (\cmd{healthy\_dog.json}) et sources.}\label{tab:rythme}\\
\toprule
Paramètre & Valeur & Source (accès) & Év.\\
\midrule\endfirsthead
\toprule
Paramètre & Valeur & Source (accès) & Év.\\
\midrule\endhead
\vp{av_erp}, \vp{av_conceal} & 0,15 s ; 0,75 & calés sur la FA canine : FC Holter médiane 129,7 [116,8--156,0], CVRR 0,309, $n = 50$ (10.3389/fcvm.2026.1890025), sous la contrainte d'une conduction 1:1 jusqu'à $\approx 225$/min (FC max Holter des sains $204 \pm 25$, 10.1371/journal.pone.0319886). Modèle : 129,6/min, CVRR 0,291 & C\\
\vp{av_symp} & 0,5 & PMC7872775 (\acc{lu}) ; Rigel 1990 (\acc{résumé}) & C\\
\vp{vent_escape_rate} & 42 /min & $42 \pm 11$, $n = 97$ (10.1111/jvim.12512) & B\\
\vp{k_atrial_standstill} & 8,5 mmol/L & arrêt atrial à 8,5--10 avec conduction sino-ventriculaire, Hariman 1983 (\acc{résumé}) & C\\
\vp{k_hp_block} & 10 mmol/L & élargissement du QRS au-delà de 10,1, Tag \& Day 2008 (\acc{moteur}) & C\\
\vp{hv} & 0,035 s & placeholder & D\\
\vp{av_beta}, \vp{av_tau}, \vp{av_vagal}, \vp{av_k_gain} & 0,1 s ; 0,15 s ; 0,5 ; 0,1 & placeholders & D\\
\vp{af_cycle}, \vp{af_cv} & 0,12 s ; 0,3 & placeholders & D\\
\vp{hp_k_slope} & 0,5 & placeholder & D\\
\vp{junct_rate} & 47 /min & au tonus de repos ; borne basse : $37 \pm 9$/min sous cholinergique local (physostigmine dans l'artère du nœud AV, $n = 10$ ; Rosenshtraukh 1986, \acc{résumé}) ; borne haute : au-dessus de l'échappement ventriculaire (42/min). Valeur retenue dans cet intervalle, sans source directe au tonus de repos & D\\
\vp{esc_sup_gain} & 1,8 & pause de 4,7 $\pm$ 1,1~s après entraînement d'un rythme jonctionnel à 2,5 fois sa fréquence, chien vigile, $n = 15$ (Tanahashi 1979, \acc{résumé}) ; pauses Holter sans échappement : 2,7~s (sains), 4,1~s (tonus vagal élevé), 6,9~s (maladie du sinus) (Flanders 2024, \acc{lu}) & C\\
\vp{esc_sup_tau} & 60 s & suppression maximale après $\approx 3$~min d'entraînement (Krellenstein 1978, cœur isolé, \acc{résumé}) ; retour en 66--100 battements jonctionnels (Neely 1985, \acc{résumé}) & D\\
\vp{esc_a_s}, \vp{esc_a_v} & 2,0 ; 0,72 & atropine et practolol sur la pause d'échappement (Tanahashi 1979, \acc{résumé}) ; voir section~\ref{sec:auto} & C\\
\vp{esc_k_gain} & 0,1 & placeholder & D\\
\vp{vpc_k_gain}, \vp{k_vpc} & 1,0 ; 3 mmol/L & placeholders & D\\
\vp{sa_block}, \vp{af}, \vp{vpc_rate}, \vp{hp_conduction} & 0 ; 0 ; 0 ; 1 & état sain (modifiés par les maladies) & --\\
\bottomrule
\end{tabularx}

\limite{Seuils de K$^+$ probablement trop hauts : un cas d'Addison (K 7,9 ; Pașca 2025, \acc{lu}) montre un BAV complet, et un essai de perfusion de KCl (Porter 2001, \acc{résumé}) une perte des ondes P vers 7,5 ; sources jugées insuffisantes pour recaler. Pas de bradycardie hyperkaliémique sous 8,5~mmol/L (FC de l'Addison 0,8 : 141/min). Pause compensatrice des ESV rarement complète (ERP nodal court, calé sur la FA). Depuis \vp{av_symp}, la FA non traitée accélère ($s = 0{,}8$ : 153 $\to$ 177/min). Les échappements suivent le tonus autonome, mais leur fréquence au repos repose sur des données expérimentales et cliniques peu nombreuses (cas isolés de Santilli 2019 : 20 à 50/min). Pas de mécanisme de réentrée, de tachycardie ventriculaire soutenue ni de voie accessoire.}

%=====================================================================
\section{Échanges gazeux et acide-base}\label{sec:gaz}
%=====================================================================

\textbf{Rôle.} PaCO$_2$ (contrôle ventilatoire), PaO$_2$ (oxygénation), pH et HCO$_3^-$ (Stewart), transport de l'O$_2$ (Hill, Bohr), sang veineux mêlé (\cmd{gas.rs}).

\subsection{Acide-base de Stewart}
\begin{gather}
\text{SID} = \vp{na} + \vp{k} - \vp{cl} - \text{lactate} - \vp{xa},\qquad [\text{H}^+] = 10^3\cdot 10^{-\text{pH}},\\
\text{HCO}_3^- = 0{,}0307\,\text{PCO}_2\,10^{\text{pH} - 6{,}1},\\
A_\text{alb} = \text{alb}_{g/L}\left(0{,}28 + \frac{0{,}469}{1 + 10^{7{,}77 - \text{pH}}}\right),\qquad A_\text{phos} = \vp{phos}\,(0{,}309\,\text{pH} - 0{,}469),\\
\text{SID} + [\text{H}^+] - \text{HCO}_3^- - A_\text{alb} - A_\text{phos} = 0,
\end{gather}
résolu en pH par Brent sur $[5{,}5;\,10{,}5]$ ($\text{alb}_{g/L} = 10\,\vp{alb}$). Excès de base (van Slyke simplifié, Hb en mmol/L $= 0{,}6206\,\vp{hb}$) : $\text{BE} = (1 - 0{,}014\,\text{Hb})\bigl(\text{HCO}_3^- - 24{,}4 + (2{,}3\,\text{Hb} + 7{,}7)(\text{pH} - 7{,}4)\bigr)$. Trou anionique de sortie : $\text{Na} + \text{K} - \text{Cl} - \text{HCO}_3^-$.

\subsection{Ventilation et oxygénation}
PaCO$_2$ résout, par Brent sur $[5;\,250]$~mmHg :
\begin{gather}
\text{PCO}_2 = 0{,}863\,\frac{\vp{vco2}\cdot\vp{bw}}{\text{VA}},\\
\text{VA} = \vp{va0}\,\vp{bw}\cdot\vp{vent_drive}\cdot\max\bigl(0{,}12;\,1 + 0{,}08(\text{PCO}_2 - 37)\bigr)\,e^{-2{,}5(\text{pH} - 7{,}4)}\,\bigl(1 + 0{,}03\max(0, 70 - \text{PaO}_2)\bigr),
\end{gather}
le pH venant de Stewart à chaque essai ; trois itérations externes sur PaO$_2$ (départ 90~mmHg) :
\begin{equation}
\text{PAO}_2 = \vp{fio2}(\vp{pb} - 47) - \text{PCO}_2/\vp{rq},\quad
\text{A-a} = \vp{aagrad} + \vp{edema_k}\max(0, \text{POG} - 18),\quad
\text{PaO}_2 = \max(15, \text{PAO}_2 - \text{A-a}).
\end{equation}
Fréquence respiratoire de sortie : $\text{FR} = 20\sqrt{\text{VA}/(\vp{va0}\,\vp{bw})}$.

\subsection{Transport de l'O$_2$}
\begin{gather}
S = \frac{P^n}{P^n + P_{50}^n},\quad n = \vp{hill_n},\quad P_{50} = \vp{p50}\cdot 10^{0{,}37(7{,}4 - \text{pH})},\\
C_aO_2 = 1{,}34\,\vp{hb}\,S + 0{,}003\,\text{PaO}_2\quad(\text{mL/dL}).
\end{gather}
\textbf{Sang veineux mêlé} : $C_vO_2 = C_aO_2 - \text{VO}_2/(\text{DC}\cdot 10)$ ; PvO$_2$ par inversion de la courbe sur $[1;\,600]$~mmHg ; $\text{PvCO}_2 = \text{PaCO}_2 + 5{,}5\,(0{,}125\,\vp{bw})/\max(\text{DC}, 0{,}2)$ ; pH veineux par Stewart avec PvCO$_2$. Particularité conservée de la référence : un contenu positif inférieur au contenu à 1~mmHg renvoie 600 et non 1.

\sources{Albumine canine : $A_\text{tot}$ 0,469~mmol/g, p$K_a$ 7,77 (Constable \& Stämpfli 2005, J Vet Intern Med 19:507, doi:10.1111/j.1939-1676.2005.tb02719.x ; \acc{résumé}) ; Na, K, Cl, albumine, phosphate du sain : même source. Effet Bohr $\Delta\log P_{50}/\Delta\text{pH} = -0{,}370$ et P50 $30{,}0 \pm 1{,}3$~mmHg (Cambier 2004, Res Vet Sci 77:83, doi:10.1016/j.rvsc.2004.02.004 ; \acc{résumé}). Gaz artériels et veineux mêlés du sain (cibles) : Haskins 2005 (\acc{résumé}). Coefficients du contrôle ventilatoire (0,08 ; 2,5 ; 0,03 ; 0,12), gradient A-a de base, seuil d'œdème 18 et pente \vp{edema_k} : non documentés.}

\textbf{Statut.} Albumine, Bohr, P50 : \evA{}. Phosphate : coefficient humain (Figge--Fencl), \evD{}. Contrôle ventilatoire : \evD{}.

\limite{Comparaison au sain (17/11/2026) : pH, PaCO$_2$, PaO$_2$, PvO$_2$ dans les cibles ; excès de base $-1{,}27$ contre $-2{,}1$ (convention de signe à confirmer) ; extraction d'O$_2$ 17,3~\% contre 20,5~\%. Fréquence respiratoire constante en ICC (23/min contre médiane 64, Yin 2026, \acc{lu}) : la tachypnée n'est pas modélisée. Seuil d'œdème 18~mmHg contre 20--24 dans les sources. Pas d'espace mort, de shunt explicite, de CO, ni de calcium.}

%=====================================================================
\section{Métabolisme et lactate}\label{sec:metab}
%=====================================================================

\textbf{Rôle.} Relier l'apport d'O$_2$ à la consommation et produire l'hyperlactatémie (\cmd{engine::metabolic\_state}).
\begin{gather}
\text{DO}_2 = \text{DC}\cdot 1000\cdot C_aO_2/100\ \ (\text{mL/min}),\qquad \text{VO}_2 = \vp{vo2}\cdot\vp{bw},\qquad \text{er} = \text{VO}_2/\text{DO}_2,\\
\text{lactate} = \vp{lact0} + \frac{9}{1 + \exp\bigl(-(\text{er} - 0{,}55)/0{,}06\bigr)}.
\end{gather}
Ordre de calcul : gaz avec \vp{lact0} $\to$ DO$_2$, er, lactate ; gaz recalculés avec ce lactate (le lactate abaisse le SID) $\to$ DO$_2$ final ; puis sang veineux mêlé. \vp{er} est celui du premier passage. Le glucose, l'ALAT et l'hématocrite ($3\,\vp{hb}$) sont des sorties directes des paramètres.

\sources{VO$_2$ cible 5,4~mL/min/kg « dérivée des REE publiées (PMC6095264) » (\cmd{targets.json}, tolérance choisie). Inflexion à er = 0,55, plateau de 9~mmol/L, pente 0,06 : non documentés (\evD{}).}

\limite{VO$_2$ fixe (pas de baisse de consommation en hypoxie, pas de fièvre ni d'exercice, sauf via les modules). Le lactate ne dépend pas de la clairance hépatique. Glucose constant, hors effet de l'Addison.}

%=====================================================================
\section{Rein}\label{sec:rein}
%=====================================================================

\textbf{Rôle.} DFG, créatinine, urée ; masse néphronique pour la clairance des médicaments et l'effet du furosémide.

\subsection{État stationnaire}
\begin{gather}
\text{DC}_\text{ref} = 3{,}29\,(\vp{bw}/20)^{0{,}75}~\text{L/min},\\
\text{perf} = \max\Bigl(0{,}25;\ \min\bigl(\mathrm{clip}(\tfrac{\text{PAM} - 40}{35}, 0, 1),\ \mathrm{clip}(\tfrac{\text{DC}/\text{DC}_\text{ref} - 0{,}5}{0{,}5}, 0, 1)\bigr)\Bigr),\\
r_\text{DFG} = \max(0{,}03;\ \vp{nephron_mass}\cdot\text{perf}),\qquad \text{DFG} = \vp{gfr0}\cdot\vp{bw}\cdot r_\text{DFG},\\
\text{créat}_{ss} = \vp{creat0}/r_\text{DFG},\qquad \text{urée}_{ss} = \vp{urea0}/r_\text{DFG}^{\,0{,}8}.
\end{gather}
La perfusion rénale est maximale au-dessus de PAM 75 et de DC $= \text{DC}_\text{ref}$ ; le plateau 0,25 traduit une azotémie prérénale, sans insuffisance rénale aiguë.

\subsection{Cinétique de la créatinine et de l'urée (18/11/2026)}
Dans les évolutions temporelles (\cmd{kinetics::lagged\_marker}, API \cmd{timecourse} et \cmd{state} avec \cmd{doses}/\cmd{at\_h}) :
\begin{equation}
V\frac{\dd C}{\dd t} = P - \text{DFG}(t)\,C \iff \frac{\dd C}{\dd t} = \frac{\text{DFG}}{V}\bigl(C_{ss}(t) - C\bigr),\qquad V = \vp{creat_vd}\cdot\vp{bw}\ \text{ou}\ \vp{urea_vd}\cdot\vp{bw}.
\end{equation}
Entre deux instants, $k = \text{DFG}\cdot 0{,}06/V$ (h$^{-1}$) et $C_{ss}$ sont pris à la moyenne des extrémités et l'équation est résolue exactement : $C \leftarrow \bar C_{ss} + (C - \bar C_{ss})e^{-\bar k\,\Delta t}$. Départ à l'équilibre du patient sans médicament ; un instant sans valeur (décès) arrête la série. L'état stationnaire (CLI \cmd{eval}, golden) reste $C_{ss}$. Non porté en Python.

\sources{$\vp{creat_vd} = 0{,}539$~L/kg (Collignon 2012, \acc{résumé}, $n = 50$) ; $\vp{urea_vd} = 0{,}66$~L/kg (Bowsher 1985, \acc{résumé}) ; aucun DOI dans le dépôt. Bornes IRIS de créatinine (IRIS 2026, \acc{lu}) pour la MRC. DFG du sain comparé à Bexfield 2008 (iohexol, 1,29--3,50~mL/min/kg, \acc{résumé}), Kelley 2022 et Kelly 2016 (\acc{lu}). Fonction de perfusion (seuils 40 et 35~mmHg, 0,5, plateau 0,25, exposant 0,8 de l'urée) : non documentée (\evD{}).}

\limite{Pas de débit urinaire de base, pas de bilan de Na$^+$, K$^+$, Cl$^-$, H$^+$ ni d'eau : en MRC, les électrolytes sont des décalages fixés par la sévérité (section~\ref{sec:ckd}). Pas de SDMA ni d'UP/C. La production de créatinine est constante (pas de masse musculaire).}

%=====================================================================
\section{Foie}\label{sec:foie}
%=====================================================================

\textbf{Rôle.} Très réduit : \vp{hepatic_fn} module la clairance des médicaments par $\vp{hepatic_fn}^{\vp{hep_exp}}$ (section~\ref{sec:pk}) ; \vp{alt} (45~U/L) est renvoyée telle quelle. Aucun module de maladie ne modifie ces paramètres. Aucune source : bloc de structure, sans contenu physiologique propre (pas de glycogénolyse, de clairance du lactate, de synthèse d'albumine).

%=====================================================================
\section{Pharmacocinétique et pharmacodynamie}\label{sec:pk}
%=====================================================================

\subsection{Pharmacocinétique commune}
Deux compartiments et un dépôt digestif par médicament (\cmd{pkpd.rs}) :
\begin{gather}
\text{CL} = \vp{cl}\cdot 10^{-3}\cdot 60\cdot 20\,(\vp{bw}/20)^{0{,}75}\cdot\vp{nephron_mass}^{\vp{renal_exp}}\cdot\vp{hepatic_fn}^{\vp{hep_exp}}~\text{(L/h)},\\
V_c = \vp{vc}\cdot\vp{bw},\quad V_p = \vp{vp}\cdot\vp{bw},\quad Q = \vp{q}\cdot 10^{-3}\cdot 60\cdot\vp{bw},\\
\dot A_g = -k_a A_g,\qquad
\dot A_c = k_a A_g - \frac{\text{CL}}{V_c}A_c - \frac{Q}{V_c}A_c + \frac{Q}{V_p}A_p,\qquad
\dot A_p = \frac{Q}{V_c}A_c - \frac{Q}{V_p}A_p,
\end{gather}
avec \vp{cl} et \vp{q} en mL/min/kg (chien de 20~kg), $C = A_c/V_c$ (mg/L). Dose IV : $A_c \mathrel{+}= \text{dose}\cdot\vp{bw}$ ; dose orale : $A_g \mathrel{+}= \text{dose}\cdot\vp{bw}\cdot F$.

\subsection{Pharmacodynamie et toxicité}
\begin{equation}
E = \frac{E_\text{max}\,C^h}{\text{EC}_{50}^h + C^h},\qquad T = \frac{E_\text{max}^\text{tox}\,C^{h_t}}{(\text{EC}_{50}^\text{tox})^{h_t} + C^{h_t}}.
\end{equation}
\cmd{apply\_drug\_effects} applique, médicament par médicament dans l'ordre alphabétique des identifiants, les effets JSON avec $E$ (et $D$, $U$ pour le diurétique), puis les termes toxiques avec $T$. L'effet thérapeutique reste plafonné (occupation des récepteurs) ; le terme toxique prolonge la relation au-delà, par un autre mécanisme. Les recouvrements entre médicaments ne sont que multiplicatifs : l'ordre ne change le résultat qu'à l'arrondi près.

{\small\setlength{\tabcolsep}{3pt}\begin{tabularx}{\linewidth}{@{}P{2.1cm}*{7}{>{\centering\arraybackslash}X}P{2.2cm}P{0.5cm}@{}}
\caption{Constantes PK et PD des modules (\cmd{modules/drugs/*.json}).}\label{tab:pk}\\
\toprule
Module & $V_c$ L/kg & $V_p$ L/kg & CL mL/min/kg & $Q$ mL/min/kg & $k_a$ /h & $F$ & renal / hep & EC$_{50}$ mg/L ; $h$ & Év.\\
\midrule\endfirsthead
\toprule
Module & $V_c$ L/kg & $V_p$ L/kg & CL mL/min/kg & $Q$ mL/min/kg & $k_a$ /h & $F$ & renal / hep & EC$_{50}$ mg/L ; $h$ & Év.\\
\midrule\endhead
furosémide 1.2.0 & 0,2073 & 0,4602 & 10,724 & 3,341 & 0,724 & 0,55 & 0,7 / 0 & 0,2412 ; 3 & B\\
pimobendane 1.0.0 & 2,0 & 1,1 & 27,9 & 16,7 & 1,08 & 0,5 & 0 / 1 & 0,03 ; 1,2 & C\\
atropine 1.0.0 & 0,7 & 1,0 & 8,0 & 10,0 & 2,0 & 0,8 & 0,3 / 0,5 & 0,004 ; 1,5 & D\\
bêtabloquant 1.1.0 & 1,0 & 1,2 & 3,0 & 5,0 & 0,5 & 0,6 & 0,9 / 0 & 0,15 ; 1,2 ($E_\text{max}$ 0,85) & D\\
opioïde 1.1.0 & 3,0 & 4,0 & 25,0 & 30,0 & 1,0 & 0,5 & 0 / 1 & 0,015 ; 1,3 & C\\
\bottomrule
\end{tabularx}}

\subsection{Furosémide (module 1.3.0, evidence B)}\label{sec:furo}
\textbf{PK.} Calée sur Koh 2021 (2~mg/kg IV et comprimé, 5 beagles d'environ 11,7~kg ; \acc{lu}, PMC8420016) : CL 735,7~mL/h/kg ramenée au chien de 20~kg par $\vp{bw}^{0{,}75}$ (10,724~mL/min/kg), $V_{ss}$ 0,6675~L/kg, $C_0$ 9,65~mg/L, $t_{1/2}$ 2,07~h ; $k_a$ et $F$ sur le comprimé ($C_\text{max}$ 0,61~mg/L, $F$ 0,55). \vp{renal_exp} = 0,7 : source non documentée.

\textbf{Effet piloté par l'excrétion urinaire} (\vp{drive}: \vp{urinary_excretion}). Le furosémide agit depuis la lumière tubulaire : la variable d'effet est $X = C\cdot\vp{nephron_mass}$ (proxy de l'excrétion urinaire), pour la diurèse comme pour les effets JSON. Chez le chien insuffisant rénal, la relation excrétion urinaire--effet est inchangée alors que la relation plasma--effet est déplacée (Miyazaki 1990, Nihon Juigaku Zasshi 52:265, doi:10.1292/jvms1939.52.265 ; Hirai 1992 ; \acc{résumés}).

\textbf{Diurèse et état de volume.} Variables d'état $D$ (déficit, mL), $U$ (diurèse effective cumulée, mL/kg), $W$ (adaptation, mL/kg) :
\begin{gather}
u = \vp{urine_max}\cdot E(X)\quad(\text{mL/kg/min}),\qquad \vp{urine_max} = 0{,}2044,\\
\frac{\dd D}{\dd t} = \vp{f_plasma}\cdot u\cdot\vp{bw}\cdot 60 - \frac{D}{\tau_\text{refill}},\qquad
\frac{\dd U}{\dd t} = 60\,u - \frac{U}{\tau_\text{elect}},\qquad
\frac{\dd W}{\dd t} = \frac{U - W}{\tau_\text{adapt}},\\
U_\text{eff} = U - \varphi\,W,
\end{gather}
par heure, avec $\vp{f_plasma} = 0{,}3$, $\tau_\text{refill} = 4$~h, $\tau_\text{elect} = 200$~h, $\tau_\text{adapt} = 12$~h, $\varphi = 0{,}9808$. Les constantes de récupération sont celles du premier diurétique dosé. La série $u$ renvoyée par \cmd{timecourse} et utilisée par les effets est $U_\text{eff}$.

\textbf{Effets} (dans l'ordre du JSON) :
\begin{gather}
\vp{vol_shift} \mathrel{-}= D,\qquad \vp{ven_cap} \mathrel{\times}= (1 + 0{,}36\,E),\\
\vp{k} \mathrel{\times}= e^{-0{,}00585\,U_\text{eff}},\qquad \vp{cl} \mathrel{\times}= e^{-0{,}00225\,U_\text{eff}},\qquad \vp{na} \mathrel{\times}= e^{-0{,}00032\,U_\text{eff}},\\
\vp{hb} \mathrel{\times}= \frac{V_b}{V_b - D},\qquad \vp{alb} \mathrel{\times}= \frac{V_p}{V_p - D}\qquad\text{(hémoconcentration, 19/11/2026)}.
\end{gather}
$V_b = 85\,\vp{bvf}\,\vp{bw}$~mL et $V_p = V_b\,(1 - \text{Ht})$ : le déficit est prélevé sur le
plasma. Le modèle donne $+7$~\% d'albumine et $+2$~\% d'hématocrite à 2~h après 2~mg/kg IV, pour
des protéines totales mesurées à $5{,}90 \pm 0{,}36 \to 6{,}32 \pm 0{,}29$~g/dL et un hématocrite
inchangé (Lee 2021, BMC Vet Res, PMC8420016, \acc{lu} ; $n = 5$) ; hématocrite et solides totaux
en hausse pendant une perfusion de 3,3~mg/kg sur 5~h (Adin 2021, PMC8163123, \acc{lu}).
Les pertes sont proportionnelles à la concentration, donc jamais négatives (correction 1 du 18/11 : K$^+$ à 200~mg/kg passe de $-0{,}60$ à 2,09~mmol/L).

\sources{
\begin{itemize}
\item Diurèse : \vp{urine_max} et EC$_{50}$ ajustés sur les volumes urinaires de Koh 2021 par intervalle (débit de base soustrait, IV et PO ensemble) ; Hill fixé à 3 (borne ; l'ajustement libre donne $\approx 5$). L'EC$_{50}$ plasmatique de 0,24~mg/L diffère de la demi-réponse de 1,5~\textmu g/mL de Hirai 1992 (chiens anesthésiés, 5~mg/kg, \acc{résumé}).
\item Électrolytes : $a_K$, $a_{Cl}$, $a_{Na}$ et $\tau_\text{elect}$ ajustés ensemble sur Koh 2021 (K, Cl, Na sériques à 1--8~h, \acc{lu}) ; Wong 1977 (1~mg/kg IV/j $\times$ 7~j : K 4,48 $\to$ 4,09 ; doi:10.1007/BF03006233, \acc{résumé}, chiens anesthésiés) n'est plus reproduit (modèle 4,49) : il entrait en conflit avec les cohortes lues.
\item Adaptation $W$ ($\tau$ 12~h, $\varphi$ 0,98 ; aldostérone, apport, excrétion fractionnelle de K) : calée sur Adin 2020 (\acc{lu} ; stade C, 3,5~mg/kg/j : K 4,50 contre 4,70, $-4$~\%) et l'essai TEST (Chetboul 2017, PMC5697186, \acc{lu} ; K 4,50 à J84 ; Adin 2020 : PMC6979094). $\tau_\text{adapt}$ = 12~h est un placeholder D : 12 et 24~h reproduisent aussi bien les cohortes, 12~h limite le creux transitoire (K 3,5 à J2 sous 2~mg/kg PO toutes les 12~h, non documenté). Test \cmd{chronic\_furosemide\_kalemia\_matches\_cardiac\_cohorts} : 1,75~mg/kg PO toutes les 12~h pendant 28~j, K$^+$ moyen entre 92 et 99,5~\% de la base ; sans adaptation, K$^+ < 3$.
\item Veinodilatation $\vp{ven_cap} \times (1 + 0{,}36E)$, $E$ suivant l'excrétion urinaire : Bayne \& Williamson 1979 (Res Commun Chem Pathol Pharmacol 25:399, \acc{résumé} : compliance veineuse périphérique 120~\% à 10~min et 136~\% à 60~min après 2~mg/kg chez le chien aux uretères ligaturés, nulle chez le néphrectomisé). L'amplitude 0,36 est la valeur publiée ; seul \vp{mr_k} est calé sur Suzuki 2013.
\item \vp{f_plasma} et $\tau_\text{refill}$ : placeholders D.
\end{itemize}}

\limite{Diurèse orale surestimée de 26~\% (27,9 contre 22,1~mL/kg sur 8~h) ; Hill à sa borne. En MRC, diurèse effective encore $\approx +25$~\% contre le sain alors que la source décrit une baisse (correction 5 « partielle »). PAM du chien sain $-10$~\% environ à 1--2~h, alors qu'une étude ne trouve pas de baisse à 1~h après 5~mg/kg (PMC10357252, 5 chiens, \acc{résumé}). Aux faibles doses (1--2~mg/kg IV), la POG baisse plus que dans Suzuki 2011 : la veinodilatation sature dès les faibles concentrations. Ototoxicité non modélisée : la relation dose--réponse existe chez le beagle IV (baisse du potentiel d'action N1 avant le potentiel microphonique, sans lésion histologique) mais aucun seuil chiffré n'est accessible (Brown 1979--1981 et Goldman 1973 : titres seulement ; Oishi 2012, PMC3490132, \acc{lu} : effet surtout transitoire). L'alcalose métabolique hypochlorémique (HCO$_3^-$ $+6{,}6$~mEq/L en 4~j, Bosch 1977, \acc{résumé} ; Adin 2021) et les arythmies hypokaliémiques (QT allongé sous 3,3~mmol/L, Hanton 2007 ; TV et torsades sous 2~mmol/L, Weissenburger 1991, \acc{résumés}) émergent des mécanismes existants (SID de Stewart, foyers ectopiques sous \vp{k_vpc}), sans calage propre.}

\subsection{Pimobendane (module 1.1.0, evidence C)}
\begin{gather}
\vp{ees_lv} \mathrel{\times}= (1 + 0{,}35\,E),\quad \vp{ees_rv} \mathrel{\times}= (1 + 0{,}20\,E),\quad \vp{r_sys} \mathrel{\times}= (1 - 0{,}20\,E),\\
\vp{pr} \mathrel{\times}= (1 - 0{,}25\,E),\qquad \vp{av_erp} \mathrel{\times}= (1 - 0{,}25\,E),\\
\text{toxique (vasodilatation, PDE3)} : \vp{r_sys} \mathrel{\times}= (1 - 0{,}8\,T),\quad \text{EC}_{50}^\text{tox} = 0{,}3~\text{mg/L},\ h_t = 3.
\end{gather}
\sources{PK calée sur Her 2020, Front Vet Sci 7:423, doi:10.3389/fvets.2020.00423 (\acc{lu} ; bras oral 0,5~mg/kg, $n = 7$) : $C_\text{max}$ 43,8 contre $49{,}1 \pm 28{,}7$~ng/mL, $t_{1/2}$ 1,66 contre $1{,}8 \pm 0{,}8$~h, ASC 149,3 contre $148{,}4 \pm 71{,}6$ ; $F = 0{,}5$ supposée. PD : placeholder. Terme toxique : EC$_{50}^\text{tox} = 10\times$ EC$_{50}$, Hill 3, aucune source (\evD{}).
Conduction AV (ajout du 19/11/2026) : PQ $93 \pm 1{,}8 \to -11{,}8$~\% à 60~min après 0,15~mg/kg IV
(Pichayapaiboon 2021, Front Vet Sci, PMC8417876, \acc{lu} ; $n = 9$), où $E \approx 0{,}45$ dans le
modèle ; réfractarités auriculaire, ventriculaire et nodale AV abaissées, AH et PR raccourcis
(Kitzen 1988, J Pharmacol Exp Ther 244:929, \acc{résumé}) ; PR et QT raccourcis à 4 fois la dose
pendant 6 mois (Kuntz 2018, \acc{résumé}). Le même facteur est appliqué à \vp{av_erp} : valeur non
publiée dans le résumé (\evD{}).}
\limite{$T_\text{max}$ 0,95~h contre $2{,}1 \pm 0{,}9$~h. Pas de voie IV dans le module. L'hypertension rapportée en surdosage (2 chiens sur 7, 2,6--21,3~mg/kg PO, base Pet Poison Helpline, \acc{résumé}) est impossible dans le modèle ; la tachycardie (4 sur 7) et l'hypotension (2 sur 7) émergent, elles, du terme vasodilatateur et du baroréflexe. Aucune concentration plasmatique après surdosage n'est publiée. L'inotropisme réduit peu la POG en MMVD à dose unique (les études portent sur des traitements de plusieurs jours).}

\subsection{Atropine (module 1.1.0, evidence C)}\label{sec:atropine}
Recalée le 19/11/2026 sur la littérature canine. \textbf{PK} : $V$ 3,2~L/kg réparti en
$\vp{vc} = 1{,}0$ et $\vp{vp} = 2{,}2$ (partage \evD{}), CL 44~mL/min/kg,
$\vp{ka_im} = \ln 2 / 7~\text{min} = 5{,}94$/h, $\vp{f_im} = 1$ (Thiermann 1996, Arch Toxicol,
doi:10.1007/s002040050276, \acc{résumé} : 2~mg IM par auto-injecteur, $C_\text{max}$ 33~ng/mL,
$t_{1/2}$ d'absorption 7~min, d'élimination 52~min) ; cohérent avec Smallridge 1989
(doi:10.1097/00003246-198912000-00003, \acc{résumé} : 0,05~mg/kg IM, pic $20{,}6 \pm 4{,}9$~ng/mL,
$V_d$ $3{,}21 \pm 0{,}63$~L/kg). Voie orale ($\vp{ka}$, $\vp{f}$) : placeholders \evD{}.

\textbf{Effet principal, au site d'effet} ($\vp{ke0} = 20$/h, $t_{1/2}$ d'équilibre $\approx 2$~min,
\evD{}) : $\vp{vagal_block} \mathrel{+}= E$ et $\vp{pupil_block} \mathrel{+}= E$, bornés à $[0, 1]$,
EC$_{50} = 0{,}0106$~mg/L, $h = 1{,}5$.

\textbf{Termes secondaires} (tableau \vp{toxicity}, chacun avec sa source dans le module) :
\begin{gather}
\text{vagal paradoxal (plasma)} : \vp{vagal_add} \mathrel{+}= 0{,}2\,T,\quad \text{EC}_{50} = 0{,}004,\ h = 2,\\
\text{antimuscarinique périphérique (site d'effet)} : \vp{musc_block} \mathrel{+}= T,\quad \text{EC}_{50} = 0{,}0015,\ h = 1{,}5,\\
\text{tachycardie excessive (plasma)} : \vp{hr} \mathrel{\times}= (1 + 0{,}2\,T),\quad \text{EC}_{50} = 0{,}045,\ h = 3.
\end{gather}
Le paradoxe de la faible dose émerge de la \emph{dissociation temporelle} : le terme vagal agit au
site plasmatique, sans délai, tandis que le blocage passe par le compartiment d'effet. Le modèle
donne donc une bradycardie initiale puis une tachycardie, comme les cohortes.
\sources{
\begin{itemize}
\item EC$_{50}$ du blocage calée pour $+50$~\% de FC à $39{,}5 \pm 6{,}2$~\textmu g/kg IV (Hendrix \& Robinson 1997, J Vet Pharmacol Ther, doi:10.1046/j.1365-2885.1997.00084.x, \acc{résumé} ; $n = 6$), en présence du terme paradoxal.
\item Blocage complet : FC $\times 1{,}74$ dans le modèle, pour $\times 1{,}9$ mesuré après 0,06~mg/kg IV cumulés (Little 1999, Am J Physiol, doi:10.1152/ajpheart.1999.276.2.H758, \acc{résumé} ; indice vagal CIPA à $1{,}3 \pm 0{,}7$~\% de la base) ; FC intrinsèque 141--146/min (Billman 2015, PMC4451292, \acc{résumé}).
\item Effet paradoxal : ralentissement du nœud sinusal et de la conduction AV à 6,25~\textmu g/kg IV (Kantelip 1985, Br J Anaesth, doi:10.1093/bja/57.2.214, \acc{résumé} ; $n = 12$) ; potentialisation de la bradycardie ($-16$~bpm) pendant 1,4~min en IV, 4,6 en IM, 9,1 en SC après 0,02~mg/kg, avec BAV2 chez 100~\% des chiens en IV (Rishniw 1996, Am J Vet Res 57:337, \acc{résumé} ; $n = 7$) ; FC abaissée chez 5/6 par 0,01~mg/kg IV sous oxymorphone (Copland 1992, \acc{résumé}).
\item Tachycardie excessive : $+26 \pm 6$~bpm au-delà de la FC de vagotomie à 0,1~mg/kg IV (Rigel 1984, Am J Physiol, doi:10.1152/ajpheart.1984.246.2.H168, \acc{résumé}) ; décrite aussi aux doses intermédiaires, mécanisme ganglionnaire évoqué (Vicenzi 1995, \acc{résumé}).
\item Antimuscarinique périphérique, plus sensible que le nœud sinusal : motilité du grêle $-93 \pm 13{,}9$~\% à 39,5~\textmu g/kg IV (Hendrix 1997) ; arrêt antral $\geq 30$~min à 0,02~mg/kg IM (Burger 2006, \acc{résumé}) ; sécrétions réduites dès 0,3--1~\textmu g/kg IV (Daly 1982, PMC2071796, \acc{résumé}) ; SOI 47,9 $\to$ 13,2~mmHg (Strombeck 1985).
\item Mydriase : aux mêmes doses que la tachycardie (de Ponti 1993, J Pharmacol Exp Ther 264:789, \acc{résumé} ; DE$_{50}$ 87~nmol/kg) ; maximum à $30{,}0 \pm 11{,}6$~min après 0,06~mg/kg IM (Kovalcuka 2015, doi:10.1111/vop.12138, \acc{résumé}) ; 0,02~mg/kg SC sans changement significatif (Kim 2015, PMC4427750, \acc{lu}).
\end{itemize}}
\limite{Aucune PK IV chez le chien vigile n'a été trouvée : la PK vient d'une étude IM. L'amplitude
de la mydriase en mm n'est publiée dans aucun résumé accessible (le test ne vérifie donc que le
sens, le pic entre 10 et 45~min et la dépression du réflexe). Pas de donnée canine sur la
relation dose--température : seule la perte évaporative de la glande nasale latérale est modélisée
(hypothèse \evD{}). Délire anticholinergique non modélisé (section~\ref{sec:periph}). Non
modélisés également : ESV et bigéminisme sous halothane (Muir 1978), aggravation de
l'hypertension des $\alpha_2$-agonistes (Alvaides 2008 ; molécules absentes du catalogue), baisse
du seuil de fibrillation ventriculaire (Yoon 1977), bronchodilatation.}

\subsection{Bêtabloquant générique, profil aténolol (module 1.1.0, evidence D)}
$\vp{beta_block} \mathrel{+}= E$, borné à $[0, 1]$, avec $E_\text{max} = 0{,}85$. Élimination surtout rénale (\vp{renal_exp} 0,9). Effets indirects : FC, inotropisme et conduction AV via $s_\text{eff}$ (sections~\ref{sec:auto}, \ref{sec:rythme}).
\sources{Profil aténolol (choix du 18/11/2026) : $T_\text{max}$ 1--2~h, effet faible mais significatif sur la FC de repos (Waterman 2019, Am J Vet Res, PMID 30801219, \acc{moteur}). Terme toxique de stabilisation de membrane retiré : propriété du propranolol (Love 2000, J Toxicol Clin Toxicol 38:421, \acc{moteur}).}
\limite{PK et PD placeholders. Surdosage devenu bénin au repos ($\times 100$ : FC $-7$~\%, PAM 103) ; aucune donnée canine chiffrée pour trancher. Hypoglycémie non modélisée.}

\subsection{Opioïde générique, chien vigile (module 1.2.0, evidence C)}
\begin{gather}
\vp{vent_drive} \mathrel{\times}= (1 - 0{,}25\,E),\qquad \vp{vagal_add} \mathrel{+}= 0{,}27\,E,\qquad \vp{r_sys} \mathrel{\times}= (1 + 0{,}36\,E),\\
\vp{t_set} \mathrel{-}= 2{,}0\,E,\qquad \vp{pupil_ew} \mathrel{+}= E,\qquad \vp{gi_inhib} \mathrel{+}= 0{,}9\,E,\qquad \vp{sedation} \mathrel{+}= 0{,}2\,E,\\
\text{toxique (dépression respiratoire)} : \vp{vent_drive} \mathrel{\times}= (1 - 0{,}8\,T),\quad \text{EC}_{50}^\text{tox} = 1~\text{mg/L},\ h_t = 3.
\end{gather}
À effet maximal : FC $-32$~\%, PAM $+6$~\%, PaCO$_2$ $+3$~mmHg ; plateau de PaCO$_2$ $\approx 65$~mmHg sans apnée à très forte concentration.
\sources{
PK de type méthadone ($V_d \approx 7$~L/kg, CL 25~mL/min/kg ; Ingvast-Larsson 2010 : 9,2~L/kg et 27,9~mL/min/kg, \acc{résumé}). PD recalée le 18/11/2026 :
\begin{itemize}
\item FC : méthadone 1~mg/kg IV, 110 $\to$ 75/min (Mills 2025, J Vet Pharmacol Ther, doi:10.1111/jvp.13515, PMC12415812, \acc{lu}) ; $-32$ à $-46$~\% (Maiante 2009, \acc{résumé}) ; perfusion à $\approx 30$~ng/mL, 87 $\to$ 55--66 (Amon 2021, PMC7814539, \acc{lu}).
\item PAM conservée ou en hausse par vasoconstriction (RVS $+57$ à $+165$~\%, Maiante 2009 ; Garofalo 2012, \acc{résumé} ; PAS inchangée, Amon 2021).
\item PaCO$_2$ : hydromorphone 0,2~mg/kg SC, max 44,4~mmHg (Wunsch 2010, Res Vet Sci 88:154, PMC4235959, \acc{lu}) ; méthadone 1~mg/kg sans changement significatif (Maiante 2009).
\item Plateau d'effet : Arndt 1984 (\acc{résumé}).
\item Terme toxique (\evD{}, borné par la littérature) : pas de hausse de PaCO$_2$ à 1~mg/kg de méthadone (Maiante 2009 ; $C \approx 0{,}33$~mg/L), ni apnée ni PaCO$_2 > 67$~mmHg avec 125 à 3000~\textmu g/kg de fentanyl IV chez le chien vigile (Bailey 1987, Anesth Analg 66:542, \acc{résumé}).
\item EC$_{50}$ thérapeutique 0,015~mg/L : placeholder cohérent avec l'antinociception au-delà de 17~ng/mL (Amon 2021).
\item Périphérie (ajouts du 19/11/2026). Consigne thermique abaissée de 2~\degre C à effet maximal : $T$ $38{,}4 \pm 0{,}1 \to 37{,}1 \pm 0{,}4$~\degre C (chute maximale $-1{,}9$) en perfusion de méthadone 0,1~mg/kg/h (Amon 2021, \acc{lu}) ; $-0{,}7$ / $-1{,}7$ / $-2{,}3$~\degre C à 0,3 / 0,5 / 1,0~mg/kg IM (cité par Amon) ; le halètement émerge de l'écart à la consigne, mécanisme avancé par Lucas 2001 (JAVMA, doi:10.2460/javma.2001.218.884, \acc{résumé}).
\item Myosis : \vp{pupil_ew} $+ E$ (noyau d'Edinger-Westphal, Lee \& Wang 1975) ; amplitude \evD{}, les données étant discordantes (section~\ref{sec:pupille}).
\item Motilité : \vp{gi_inhib} $+ 0{,}9E$, inhibition du complexe moteur migrant complète à 500--1000~\textmu g/kg/h de morphine, résistante à la vagotomie (Sarna 1985, \acc{résumé}) ; régurgitation chez 4/7 en perfusion de méthadone (Amon 2021).
\item Sédation légère seulement chez le vigile : \vp{sedation} $+ 0{,}2E$, soit un score de 2 à 7 sur 21 (hydromorphone 0,1~mg/kg IM non plus sédative que le sérum physiologique, Hofmeister 2010, \acc{résumé} ; méthadone 0,5~mg/kg IV légère, KuKanich 2008).
\end{itemize}}
\limite{Pas d'état « anesthésié » : l'apnée décrite sous anesthésie halogénée (Garofalo 2012) n'est pas reproductible, et \vp{anesth_mac} n'agit que sur la périphérie. Dysphorie et vomissements non modélisés (section~\ref{sec:periph} ; l'opioïde générique a un profil méthadone, dont les vomissements sont rares : 0/6 à 0/8). Histaminolibération de la morphine IV rapide (histamine 0,55 $\to$ 531~ng/mL à 2~mg/kg, Robinson 1988, \acc{résumé}) non modélisée. Le point fixe de décompensation est fragile en association (maladie du sinus + opioïde massif : décès par spirale à $\times 10$ et $\times 100$, pas à $\times 30$).}

%=====================================================================
\section{Périphérie : thermorégulation, pupille, digestif, SNC}\label{sec:periph}
%=====================================================================
Bloc \cmd{periph} (ajouté le 19/11/2026, utile pour l'anesthésie). Ces quatre effecteurs lisent
les paramètres (modifiés par les médicaments et par le scénario : ambiance, anesthésie,
éclairage) et le tonus autonome du régime permanent, puis produisent des indicateurs cliniques.
Seule la température revient sur le reste du modèle. Entrées de scénario :
\vp{t_amb}, \vp{warming}, \vp{anesth_mac}, \vp{light}, \vp{excitation} ; entrées pilotées par les
médicaments : \vp{t_set}, \vp{musc_block}, \vp{pupil_block}, \vp{pupil_ew}, \vp{gi_inhib},
\vp{sedation}.

\medskip\noindent\textbf{Avertissement.} \vp{anesth_mac} n'agit que sur la périphérie : les effets
hémodynamiques et respiratoires des halogénés ne sont pas modélisés.

\subsection{Thermorégulation}\label{sec:thermo}
Bilan thermique à deux compartiments, noyau et enveloppe ($T_c$, $T_s$ en \degre C) :
\begin{gather}
f_c\,c\,\frac{\dd T_c}{\dd t} = M - E - K\,(T_c - T_s),\qquad
(1 - f_c)\,c\,\frac{\dd T_s}{\dd t} = K\,(T_c - T_s) - G\,(T_s - T_a) + W,\\
M = M_0\,(1 - 0{,}10\,I)\,Q_{10}^{(T_c - 38{,}6)/10} + M_0\,\text{frisson},\qquad
E = \big(0{,}20\,M_\text{basal} + M_0\,\text{halètement}\big)\,(1 - 0{,}27\,\vp{musc_block}),\\
K = K_0\,e^{(1 - I)\,\text{clip}(3\,(T_c - \vp{t_set}),\ \ln 0{,}2,\ \ln 3)},\qquad I = \text{clip}(\vp{anesth_mac}, 0, 1),
\end{gather}
avec $M_0 = \vp{vo2}\cdot 20{,}1/60$~W/kg, $c = 3470$~J/kg/\degre C, $f_c = 0{,}69$. Les conductances
$K_0$ (noyau $\to$ enveloppe, $\propto \vp{bw}^{-0{,}25}$) et $G$ (enveloppe $\to$ ambiance,
$\propto \vp{bw}^{-1/3}$) sont fixées par l'équilibre du chien de référence en neutralité
(38,6 et 34,6~\degre C à $T_a = 26$~\degre C). Réponses proportionnelles à l'écart à la consigne
(Hammel 1963), supprimées par l'anesthésie ($1 - I$) :
\begin{gather}
\text{frisson} = (1 - I)\cdot 3\cdot\text{clip}\!\left(\frac{\vp{t_set} - T_c - 0{,}5}{1{,}5}, 0, 1\right),\qquad
\text{halètement} = (1 - I)\cdot 4\cdot\text{clip}\!\left(\frac{T_c - \vp{t_set} - 0{,}2}{1{,}0}, 0, 1\right).
\end{gather}
Intégration RK4 au pas de 60~s ; l'équilibre initial est obtenu sur 72~h. \cmd{thermal\_course}
applique à chaque pas les effets des médicaments aux concentrations de \cmd{timecourse}.

\textbf{Retour sur le modèle} (\cmd{apply\_temperature}, appelé en tête de \cmd{evaluate}) :
\begin{gather}
\vp{hr} \mathrel{\times}= 1{,}51^{\Delta T/10},\qquad
\vp{vo2} \mathrel{\times}= (1 - 0{,}10\,I)\,1{,}7^{\Delta T/10} + \text{frisson},\qquad
\vp{p50} \mathrel{\times}= 10^{0{,}021\,\Delta T},
\end{gather}
$\Delta T = T_c - 38{,}6$ ; identité exacte à 38,6~\degre C sans anesthésie ni frisson. La CAM
effective suit $\vp{anesth_mac}\cdot e^{-0{,}086\,\Delta T}$.

\sources{
\begin{itemize}
\item $T_c$ de référence 38,6~\degre C, intervalle 37,7--39,5 (Dorn 2024, Animals 14:1970, doi:10.3390/ani14131970, PMC11240782, \acc{lu} ; $n = 9117$) : \evA{}.
\item $f_c = 0{,}69$ calé sur la chute sous anesthésie : $38{,}7 \pm 0{,}6 \to 36{,}7 \pm 1{,}3$ / $36{,}1 \pm 1{,}4$ / $35{,}8 \pm 1{,}5$~\degre C à 1 / 2 / 3~h (Redondo 2012, Vet Rec 171:374, doi:10.1136/vr.100476, \acc{résumé} ; $n = 1525$) ; profil en plateau vers 75~min (Sage 2023, \acc{résumé}, $n = 1097$) : \evC{}.
\item $Q_{10}$ de la VO$_2$ = 1,7 : 4,7 $\pm$ 0,4 $\to$ 3,2 $\pm$ 0,3~mL/kg/min de 37,7 à 30,5~\degre C (Gutierrez 1986, J Appl Physiol 60:751, \acc{résumé}) : \evC{}.
\item Frisson $\times 4$ : VO$_2$ $+300$~\% (Auld 1980, Clin Sci 58:501, \acc{résumé}) ; $\times 4{,}5$ à $-21$~\degre C (Minaire 1982, \acc{résumé}) : \evB{}.
\item Neutralité thermique 26--29~\degre C l'été, 22--26 l'hiver (Sugano 1981, Jpn J Physiol 31:465, \acc{résumé}) : \evC{}.
\item Glande nasale latérale : 19--36~\% de la hausse d'évaporation respiratoire (Blatt 1972, Science 177:804, \acc{résumé}) ; sa commande cholinergique est une hypothèse (\evD{}).
\item $Q_{10}$ de la FC = 1,51 : FC $-17$~\% de 38,5 à $\approx 34$~\degre C sous isoflurane (Yoshida 2001, Anesthesiology 94:678, \acc{résumé} ; $n = 30$), valeur dérivée : \evC{}.
\item P50 : $\Delta\log_{10}\text{P50}/\Delta T = 0{,}022$ (Reeves 1982, \acc{résumé}) et 0,020 (Cambier 2004, \acc{résumé}) : \evB{}.
\item CAM : 1,20 $\pm$ 0,17 contre 0,89 $\pm$ 0,18~\% à 34,5~\degre C, $-26$~\% (Wilson 2006, JAVMA 228:1042, \acc{résumé}) : \evC{}.
\item Réchauffement actif (\vp{warming}, W/kg) : circuit chauffé et couverture à eau, 37,7 $\pm$ 0,6 contre 36,6 $\pm$ 1,0~\degre C à l'extubation (Jones 2023, JAVMA 261:1181, \acc{résumé}).
\end{itemize}}
\limite{Aucune valeur canine publiée n'a été trouvée pour la chaleur spécifique du corps
(3470~J/kg/\degre C : valeur humaine, \evD{}), la conductance sèche, les seuils de
vasoconstriction et de frisson (au repos comme sous anesthésie ; seul substitut trouvé : lapin,
Imamura 2003) et les pertes respiratoires en W/kg. $f_c$, $T_s$ neutre (34,6), la part
respiratoire (0,20), les gains vasomoteurs et les plages de recrutement sont donc des
placeholders \evD{} calés sur la seule cinétique de refroidissement. Le modèle de Potter 2020
(J Therm Biol 91:102651, payant) contient probablement ces constantes.}

\subsection{Pupille}\label{sec:pupille}
Innervation réciproque (Yoshitomi 1986, IOVS 27:83, \acc{résumé}) : dilatateur ($\alpha_1$, tonus
sympathique total $s$) contre sphincter (muscarinique M3, commande d'Edinger-Westphal modulée par
l'éclairement) :
\begin{gather}
D = D_\text{min} + (D_\text{max} - D_\text{min})\,\frac{d}{d + \sigma},\qquad
d = d_0\,\frac{\max(s, 0)}{0{,}2},\\
\sigma = (1 + \vp{pupil_ew})\,\big(\sigma_\text{obsc} + (1 - \sigma_\text{obsc})\,\vp{light}\big)\,(1 - \vp{pupil_block}),
\end{gather}
$D_\text{min} = 1{,}0$ et $D_\text{max} = 11{,}0$~mm (myosis et mydriase maximaux). Les gains $d_0$,
$\sigma_\text{obsc}$ et celui du flash sont fixés exactement par trois diamètres mesurés :
8,53~mm en lumière ambiante, 10,06 à l'obscurité, 6,57 au creux du réflexe photomoteur. Le modèle
reproduit alors une amplitude de réflexe de 23~\%, pour $23{,}5 \pm 7{,}0$~\% mesurés.
\sources{Mills 2022, Front Vet Sci, doi:10.3389/fvets.2022.1020710, PMC9643214 (\acc{lu} ; 126 yeux,
209--236~lux : MAX $8{,}53 \pm 1{,}30$, MIN $6{,}57 \pm 1{,}41$, constriction $23{,}5 \pm 7{,}0$~\%) : \evA{}.
Obscurité $10{,}06 \pm 0{,}54$ (Whiting 2013, cité par Mills 2022) et extrêmes pharmacologiques
$11{,}02 \pm 0{,}83$ / $1{,}02 \pm 0{,}32$ (Weisse 1975, doi:10.1007/BF00501796, \acc{résumé} ;
142 beagles) : \evB{}. Myosis opioïde par excitation du noyau d'Edinger-Westphal, aboli par section
du nerf III (Lee \& Wang 1975, J Pharmacol Exp Ther 192:415, \acc{résumé}).}
\limite{Statique : la cinétique du réflexe (latence $0{,}22 \pm 0{,}05$~s, vitesses) n'est pas
simulée. Aucune relation diamètre--concentration d'isoflurane publiée chez le chien ; la position
du globe et les réflexes palpébral et cornéen, qui servent en pratique à juger la profondeur
d'anesthésie, ne sont pas modélisés. Les amplitudes de \vp{pupil_ew} et \vp{pupil_block} sont des
\evD{} : les données opioïdes sont discordantes (morphine : myosis marqué ; méthadone 0,3~mg/kg :
aucun changement, Raušer 2022, \acc{résumé}).}

\subsection{Digestif et sécrétions}
\begin{gather}
\text{motilité} = (1 - \vp{musc_block})(1 - \vp{gi_inhib})\,(1 - 0{,}727\,\min(\text{CAM}_\text{eff}, 1)),\\
\vp{les} = 47{,}9\,(1 - 0{,}724\,\vp{musc_block})~\text{mmHg},\qquad \text{salivation} = 1 - \vp{musc_block}.
\end{gather}
\sources{Pression du sphincter œsophagien inférieur $47{,}9 \pm 1{,}2$~mmHg, $13{,}2 \pm 2{,}0$ sous
atropine (Strombeck 1985, PMID 4014848, \acc{résumé}) : \evB{}. Force de contraction gastrique
11~(8--20) $\to$ 3~(1--10)~mmHg sous sévoflurane, retour en 12--15~h (Boscan 2014,
doi:10.1111/vaa.12093, \acc{résumé} ; $n = 12$) : \evC{}. Arrêt de la motilité antrale $\geq 30$~min
à 0,02~mg/kg d'atropine IM, $> 3$~h à 0,04 (Burger 2006, doi:10.1111/j.1439-0442.2006.00787.x,
\acc{résumé}) ; motilité du grêle $-93 \pm 13{,}9$~\% à 39,5~\textmu g/kg IV (Hendrix 1997,
\acc{résumé}). Inhibition opioïde du complexe moteur migrant, dose-dépendante, complète à
500--1000~\textmu g/kg/h de morphine et résistante à la vagotomie (Sarna 1985, PMID 2993596,
\acc{résumé}) : \evC{}.}
\limite{Indices relatifs, sans débit salivaire absolu : aucune valeur canine en mL/min chez le
vigile, et aucune courbe dose--réduction de la salive sous atropine ou glycopyrrolate aux doses
cliniques n'a été trouvée. Vomissements et régurgitation non modélisés (fréquences publiées
seulement par protocole : morphine 70--76~\%, hydromorphone 0,2~mg/kg IM 87~\%, méthadone
0/6--0/8). Iléus postopératoire et vidange gastrique non simulés.}

\subsection{Système nerveux central}
\begin{gather}
\text{drive} = \vp{sedation} + \left(\frac{\text{CAM}_\text{eff}}{0{,}4}\right)^{4},\qquad
\text{score} = 1 + 20\,\frac{\text{drive}}{1 + \text{drive}}.
\end{gather}
Le BIS est interpolé dans la table mesurée et vaut \cmd{None} hors du domaine 0,8--2,0~CAM.
\sources{Échelle de sédation de Grint, 7 items, 0--21 (Wagner 2017, BMC Vet Res,
doi:10.1186/s12917-017-1027-2, PMC5395740, \acc{lu} ; $\alpha$ de Cronbach 0,89, base médiane 1,
dexmédétomidine + hydromorphone 11~[1--18], acépromazine + hydromorphone 5~[0--15], perte du tonus
de la mâchoire à 17--18/21) : \evB{}. BIS selon la CAM d'isoflurane : $65 \pm 8$ / $60 \pm 7$ /
$52 \pm 3$ / $31 \pm 28$ à 0,8 / 1,0 / 1,5 / 2,0~CAM (Greene 2003, Am J Vet Res,
doi:10.2460/ajvr.2003.64.316, \acc{résumé} ; 6 chiens curarisés) : \evC{}.}
\limite{L'exposant 4 et la CAM de demi-effet (0,4) sont des \evD{} : aucune relation
concentration--score publiée chez le chien. \vp{excitation} reste une entrée de scénario : le
délire anticholinergique (un seul cas canin, Fedotova 2025, avec lidocaïne associée et des seuils
de source secondaire contradictoires) et la dysphorie des opioïdes (présente à 1~mg/kg IV de
méthadone, Garofalo 2012 ; absente à 0,5, KuKanich 2008, \acc{résumés}) ne sont pas modélisés,
faute d'échelle et de relation concentration--effet. Convulsions non modélisées (la seule donnée
canine chiffrée porte sur la lidocaïne : $20{,}8 \pm 4{,}0$~mg/kg, Feldman 1989).}

%=====================================================================
\section{Maladies}\label{sec:maladies}
%=====================================================================

\textbf{Fonctionnement commun} (\cmd{case.rs}). Un cas tire un stade selon \vp{prior}, puis $s$ uniforme dans les bornes du stade (ou $s$ imposé, le stade se déduit des bornes). Les \vp{draws} sont tirés une fois par cas ($\mathrm{lognormal\_mean1}$ : $\exp(\mathcal N(-\sigma^2/2, \sigma^2))$ ; Bernoulli), les \vp{let} évalués dans l'ordre, puis les effets appliqués dans l'ordre (\vp{mul} : $p \leftarrow p\cdot\text{expr}$ ; \vp{add} : $p \leftarrow p + \text{expr}$). Une formule de maladie ne lit pas les paramètres du patient. Mode déterministe : lognormale $= 1$, Bernoulli $= 0$. Les signes cliniques (souffle, œdème, azotémie\ldots) émergent du moteur ; les \vp{targets} servent à la vérification, pas au calcul.

\subsection{MMVD (\cmd{mmvd.json} 1.0.0, evidence C)}\label{sec:mmvd}
Stades ACVIM : B1 $[0{,}05; 0{,}3]$, B2 $[0{,}3; 0{,}6]$, C $[0{,}6; 0{,}85]$, D $[0{,}85; 1]$ (priors 0,35 ; 0,30 ; 0,25 ; 0,10). Tirages $z_g$, $z_{cpv}$, $z_{bvf}$ ($\sigma$ 0,25). Rampes : $\ell_\text{late} = \max(0, (s - 0{,}8)/0{,}2)$, $\ell_\text{ret} = \max(0, (s - 0{,}55)/0{,}45)$, $\ell_\text{aut} = \max(0, (s - 0{,}4)/0{,}6)$.
\begin{gather}
\vp{g_mr} \mathrel{+}= 0{,}9458\,s^{1{,}111}\,z_g,\qquad \vp{c_pv} \mathrel{\times}= 1 + 0{,}116\,s\,z_{cpv},\qquad \vp{v0_lv} \mathrel{\times}= 1 + 1{,}115\,s,\\
\vp{b_lv} \mathrel{\times}= 1 - 0{,}2\,s,\qquad \vp{ees_lv} \mathrel{\times}= 1 - 0{,}35\,\ell_\text{late},\qquad \vp{bvf} \mathrel{\times}= 1 + 0{,}05\,s + 0{,}30\,\ell_\text{ret}\,z_{bvf},\\
\vp{s0} \mathrel{+}= 0{,}5\,\ell_\text{aut},\qquad \vp{v0} \mathrel{+}= -0{,}5\,\ell_\text{aut},\qquad \vp{baro_sens} \mathrel{\times}= 1 - 0{,}5\,\ell_\text{late}.
\end{gather}
Régurgitation croissante, dilatation du VG, rétention hydrosodée tardive, déséquilibre autonome ; le souffle, l'œdème, l'hypoxémie et la tachycardie émergent. Avec l'orifice dynamique, la conductance effective vaut $\vp{g_mr}\,(\text{VTD}/\vp{mr_vref})^{2{,}5}$ : les coefficients 0,9458 et 1,111 ont été recalés le 18/11/2026 pour conserver la fraction régurgitée et la POG de chaque stade (RF B1/B2/C/D 0,11/0,35/0,58/0,76 contre 0,04/0,36/0,64/0,75 avant ; POG 5,2/7,1/12,0/21,0 contre 4,9/7,3/13,5/20,0~mmHg).
\sources{Calage initial (\cmd{fit\_mmvd.py}) sur des cibles de fraction de régurgitation et de POG par stade : RF 17--81~\% en IM expérimentale (Oyama 2004, J Vet Intern Med 18:667, doi:10.1111/j.1939-1676.2004.tb02604.x, \acc{résumé}) ; œdème si POG $> \approx 20$~mmHg (MSD Vet Manual) ; FC 130 et 160/min aux stades C et D : choix clinique non sourcé. Cibles de vérification du module : RF B1 $\le 20$~\% ; POG C 14--24 ; POG D 20--32 ; FC D 140--180 (« choix clinique non sourcé »). Stades : ACVIM 2019 (Keene, doi:10.1111/jvim.15488, \acc{lu}). Comparaisons (17/11) : RF B1 7,0~\% contre 13,5~\% (IQR 1,9--28,3), B2 38,3 contre 53,3 (38,7--65,8), C 63,9 contre 60,2 (Kawai 2026, \acc{lu}) ; FC et PAS en ICC conformes à Yin 2026 (\acc{lu}).}
\textbf{Recalage 1.2.0 (19/11/2026), après le solveur du baroréflexe et le recalage de la systole.} Les coefficients du déséquilibre autonome ($\vp{s0} \mathrel{+}= 0{,}65\,\ell_\text{aut}$, $\vp{v0} \mathrel{-}= 0{,}65\,\ell_\text{aut}$ ; 0,5 auparavant) et de la rétention tardive ($0{,}37\,\ell_\text{ret}\,z_{bvf}$ ; 0,30) sont recalés. La rétention est plafonnée par l'effet du furosémide sur la POG (Suzuki 2013 : $-27$~\% ; test à $-18$~\%), qui faiblit quand la volémie augmente : à 0,38, l'effet tombe à $-17{,}8$~\%.
\begin{itemize}
\item \textbf{FC.} Les FC publiées sont mesurées en consultation : les témoins y sont à $115 \pm 22$/min (Baisan 2023, PMC10210847, \acc{lu}), alors que le modèle représente le repos (87/min, Haskins 2005). La comparaison porte donc sur l'élévation \emph{relative}. Repères (\acc{lu}) :
  \begin{itemize}
  \item B1 120 (104--132) et B2 128 (117--140), $n = 908$ et 337 (Wilshaw 2021, PMC7995403) ;
  \item B2 $142 \pm 25$ et œdème aigu $160 \pm 21$ (Baisan 2023) ;
  \item apparition de l'ICC 150 (140--162), soit $+20$~\% en dix mois (Boswood 2018, PMC5787203 ; Boswood 2020) ;
  \item récidive d'ICC, C avancé ou D : $165 \pm 31$ (Park 2026, PMC13015027) ;
  \item aucune FC publiée au stade D dans les sources lues.
  \end{itemize}
  Modèle, en rapport à la FC du chien sain : B1 $\times 0{,}99$, B2 $\times 1{,}05$ (cible 1,11--1,23 : trop bas, le déséquilibre ne démarre qu'à $s = 0{,}4$), C $\times 1{,}28$ (cible 1,30--1,39), D $\times 1{,}62$ (cible $1{,}43 \pm 0{,}27$).
\item \textbf{Cible de FC du stade D.} L'ancienne cible absolue 140--180 (« choix clinique non sourcé ») est remplacée par 100--150/min au repos, soit $87 \times (1{,}43 \pm 0{,}27)$ ; le modèle donne 141.
\item \textbf{Récapitulatif par stade} ($s$ = 0,175/0,45/0,725/0,925) : RF 0,10/0,30/0,56/0,74 ; POG 6,0/7,3/13,9/26,1~mmHg (cibles de plausibilité C : 14--24, D : 20--32 ; le milieu du stade C est à 0,1 sous la borne, au niveau de la médiane d'Allen 2021 en ICC, 14) ; FC 86/91/112/141 ; PAM 103/105/117/120.
\end{itemize}

\textbf{Congestion par stade (18/11/2026).} La rétention hydrosodée (coefficient de $\ell_\text{ret}$) est portée de 0,149 à 0,30 pour que le stade C soit congestif : POG 9,8 à 22,8~mmHg sur le stade C (15,2 au milieu, contre 12,1 avant), jusqu'à 37~mmHg au stade D ; stades B inchangés (POG 7,1 à $s = 0{,}45$). Repères : POG invasive chez des chiens MMVD en ICC (stades C/D) : médiane 14~mmHg (8--32), Allen 2021, JAVMA 258:638, doi:10.2460/javma.258.6.638 (\acc{résumé}) ; IM expérimentale de stade B2 : 14,6~$\pm$~3,8~mmHg (Suzuki 2012, PMC3489586, \acc{lu}) ; œdème expérimental aigu au-delà de 20--25~mmHg (Nakata 1986 ; Magno \& Szidon 1976 ; \acc{résumé}). Aucune série invasive par stade ACVIM n'a été trouvée chez le chien MMVD spontané. La fuite mitrale est normalisée au chien de 20~kg (section~\ref{sec:valves}) : au milieu du stade C, la POG ne dépend plus guère du poids (10 à 60~kg).
\limite{Indices échographiques (LA:Ao, LVIDDN) et NT-proBNP non calculés. Pas de tachypnée en ICC. VFC du stade C : voir section~\ref{sec:auto}. Les comparaisons de RF par stade (Kawai 2026) précèdent les recalages du 18/11. Au stade D, hypoxémie marquée (PaO$_2$ 54~mmHg au milieu du stade, 29 à $s = 1$).}

\subsection{Maladie rénale chronique (\cmd{ckd.json} 1.0.0, evidence B)}\label{sec:ckd}
Stades IRIS 1 à 4 : $[0{,}05; 0{,}25]$, $[0{,}25; 0{,}5]$, $[0{,}5; 0{,}75]$, $[0{,}75; 0{,}95]$. Tirages lognormaux ($\sigma$ 0,25) $z_{neph}$, $z_{xa}$, $z_{phos}$, $z_{hb}$, $z_{hyp}$ ; hypertension Bernoulli de probabilité $0{,}2 + 0{,}6\,s$.
\begin{gather}
\vp{nephron_mass} \mathrel{\times}= \max\bigl(0{,}05;\ e^{-2{,}4\,s\,z_{neph}}\bigr),\quad \vp{xa} \mathrel{+}= 6\,s^2 z_{xa},\quad \vp{phos} \mathrel{+}= 1{,}6\,s^2 z_{phos},\\
\vp{hb} \mathrel{\times}= 1 - 0{,}35\,s^{1{,}5} z_{hb},\quad \vp{alb} \mathrel{\times}= 1 - 0{,}1\,s,\quad \vp{k} \mathrel{+}= 2\max(0, (s - 0{,}85)/0{,}15),\\
\vp{r_sys} \mathrel{\times}= 1 + 0{,}3\,s\,z_{hyp}\quad\text{si hypertension}.
\end{gather}
\sources{Bornes IRIS de créatinine (stade 1 $< 125$, 2 : 125--250, 3 : 251--440, 4 $> 440$~\textmu mol/L ; IRIS 2026, \acc{lu}) : seul calage publié du module (\evB{}). Les autres coefficients (acidose, phosphate, anémie, albumine, K$^+$, hypertension) : non documentés.}
\limite{(Vérification du 17/11.) Le haut de la plage IRIS1 tombe déjà en stade 2 (créatinine 140 à $s = 0{,}25$). Na$^+$ et Cl$^-$ constants ; K$^+$ constant jusqu'à $s = 0{,}85$ (alors que les sources rapportent 4,8 contre 4,3 dès le stade 2) ; acidose (HCO$_3^- < 22$) seulement au stade 3 contre 67~\% des chiens dès le stade 2 (Lippi 2023, \acc{lu}) ; hématocrite trop haute de 4 à 5 points (Lippi 2021) ; calcium absent ; hypertension tirée au hasard.}

\subsection{Hypoadrénocorticisme -- Addison (\cmd{addison.json} 1.0.0, evidence C)}
Stades discret $[0{,}1; 0{,}4]$, modéré $[0{,}4; 0{,}7]$, crise $[0{,}7; 1]$. Tirages $z_c$ (partagé par Na et Cl), $z_k$, $z_{bvf}$ ($\sigma$ 0,2), hypoglycémie Bernoulli 0,5.
\begin{gather}
\vp{na} \mathrel{+}= -22\,s\,z_c,\quad \vp{cl} \mathrel{+}= -17\,s\,z_c,\quad \vp{k} \mathrel{+}= 3{,}4\,s\,z_k,\\
\vp{bvf} \mathrel{\times}= 1 - 0{,}22\,s\,z_{bvf},\quad \vp{r_sys} \mathrel{\times}= 1 - 0{,}25\,s,\quad \vp{glucose} \mathrel{+}= -2\,s\cdot\text{hypo}.
\end{gather}
\sources{Na:K $< 27$ évocateur (MSD Vet Manual ; cible du stade modéré) ; K 6--9 en crise (« cas rapportés, ex. K 8,3 », plausibilité). Comparaisons (17/11) : Na:K 23,1 à $s = 0{,}55$ sous les seuils d'Adler 2007 (\acc{résumé}) ; Na$^+$ 131,3 et K$^+$ 6,99 en crise contre médianes 137 et 7,3 (Baumstark 2014, \acc{lu}) ; créatinine 1,58~mg/dL contre 1,7 (Hall 2023, \acc{lu}).}
\limite{Crise très sensible au déficit de volume (solutions non monotones). Depuis la suppression de \vp{kfac}, tachycardie compensatrice forte (141/min à $s = 0{,}8$) et aucune bradycardie sous 8,5~mmol/L (9 à 21~\% de chiens bradycardes en crise, Chalifoux 2023).}

\subsection{Pneumonie (\cmd{pneumonia.json} 1.0.0, evidence D)}
$\vp{aagrad} \mathrel{+}= 70\,s^{1{,}3} z_{aa}$ ($\sigma$ 0,25) ; \vp{vco2} et \vp{vo2} $\times (1 + 0{,}1\,s)$ ; $\vp{lact0} \mathrel{+}= 0{,}6\,s$. Aucune référence (\vp{refs} vide) : module entièrement placeholder. Pas de fièvre, d'inflammation ni de tachypnée.

\subsection{Fibrillation atriale (\cmd{atrial\_fibrillation.json} 1.0.0, evidence C)}
$\vp{af} \mathrel{+}= 1$ ; $\vp{av_erp} \mathrel{\times}= (1{,}25 - 0{,}5\,s)\,z_{av}$ ($\sigma$ 0,1). Stades : réponse lente, modérée, rapide. Réponse ventriculaire 105, 132 et 155/min aux sévérités 0,1 ; 0,5 ; 0,9 (avant \vp{av_symp}).
\sources{FC Holter médiane 129,7 [116,8--156,0], CVRR 0,309 (Gunasekaran 2026, Front Cardiovasc Med, 10.3389/fcvm.2026.1890025) ; FC $> 150$ chez 80~\% des chiens non traités (ORCA, 10.1111/jvim.16666).}
\limite{Pas de remodelage atrial ni de lien avec la MMVD ; amplitude des ondes f variable par la seule conduction (placeholder).}

\subsection{Bloc atrioventriculaire (\cmd{av\_block.json} 1.0.0, evidence C)}
\begin{gather}
\vp{pr} \mathrel{+}= 0{,}12\min(1, s/0{,}33),\qquad \vp{av_erp} \mathrel{\times}= 1 + 4\min\bigl(1, \max(0, (s - 0{,}33)/0{,}33)\bigr),\\
\vp{hp_conduction} \mathrel{\times}= \max\bigl(0;\ 1 - \max(0, s - 0{,}66)/0{,}1\bigr).
\end{gather}
Rythmes produits : 1\textsuperscript{er} degré (PR jusqu'à 0,24~s), Wenckebach vers $s \approx 0{,}5$, 2:1 vers 0,6, 3\textsuperscript{e} degré avec échappement ventriculaire à 42/min.
\sources{PR $> 0{,}13$~s pour le 1\textsuperscript{er} degré (10.3389/fvets.2025.1622850) ; échappement $42 \pm 11$/min, QRS 61~ms, $n = 97$ (10.1111/jvim.12512).}

\subsection{Extrasystoles ventriculaires (\cmd{ventricular\_premature.json} 1.0.0, evidence D)}
$\vp{vpc_rate} \mathrel{+}= 30\,s^2 z_{vpc}$ /min ($\sigma$ 0,3). Repère : moins de 20 ESV/24~h chez 86~\% des chiens sains (10.1371/journal.pone.0319886). Coefficient 30 non documenté.

\subsection{Maladie du sinus (\cmd{sick\_sinus.json} 1.0.0, evidence D)}
$\vp{sa_block} \mathrel{+}= 0{,}45\,s\,z_{sa}$ ($\sigma$ 0,2) ; $\vp{hr} \mathrel{\times}= 1 - 0{,}3\,s$ ; \vp{junct_rate} et \vp{vent_escape_rate} $\times (1 - 0{,}8\,s)$.
\sources{Plus longue pause 6,9~s (5,5--10,3) sur Holter de 24~h, $n = 21$ (10.1111/jvim.17071) ; pauses $2{,}5 \pm 1{,}3$~s à l'ECG (10.1016/j.jvc.2018.12.003). Modèle : 4,9~s sur 5~h à $s = 0{,}85$.}
\limite{Coefficients non calés ; réponse à l'atropine insuffisante (aucun stade n'atteint 135/min).}

%=====================================================================
\section{Décompensation et viabilité}\label{sec:decomp}
%=====================================================================

\textbf{Rôle.} Faire agir l'aval métabolique en retour sur la pompe et les vaisseaux, et décider de la viabilité, de façon entièrement mécanistique (pas de risque statistique).

\subsection{Point fixe de décompensation}
\cmd{evaluate} enveloppe \cmd{evaluate\_once} d'un point fixe amorti sur $f_c$ (multiplie \vp{ees_lv} et \vp{ees_rv}) et $f_r$ (multiplie \vp{r_sys}), départ $f_c = f_r = 1$ :
\begin{align}
f_\text{acid} &= \max\bigl(0{,}1;\ 1 - \vp{acid_ees}\max(0, \vp{acid_ph0} - \text{pH}_a)\bigr) &&\text{(sur } f_c\text{)},\\
f_\text{acid}^{r} &= \max\bigl(0{,}1;\ 1 - \vp{acid_rsys}\max(0, \vp{acid_ph0} - \text{pH}_a)\bigr) &&\text{(sur } f_r\text{)},\\
f_\text{hyp} &= \max\bigl(0{,}1;\ 1 - \vp{hypox_ees}\max(0, \vp{hypox_pvo2} - \text{PvO}_2)\bigr) &&\text{(sur } f_c\text{)},\\
f_\text{cor} &= \max\bigl(0{,}1;\ 1 - \vp{cpp_ees}\max(0, \vp{cpp_min} - (\text{PAD} - \text{PTDVG}))\bigr) &&\text{(sur } f_c\text{)},
\end{align}
cibles $f_c^* = f_\text{acid}f_\text{hyp}f_\text{cor}$, $f_r^* = f_\text{acid}^r$ ; $f \leftarrow \tfrac12(f + f^*)$ jusqu'à $|f^* - f| < 10^{-3}$, au plus 30 passes. Faiblesse hypokaliémique, appliquée directement : $\vp{vent_drive} \mathrel{\times}= \max\bigl(0{,}1;\ 1 - \vp{k_weak_vent}\max(0, \vp{k_weak} - K)\bigr)$. Sans couplage actif, une seule passe, identique à \cmd{evaluate\_once}.

\subsection{Critères de décès (\cmd{Error::NonViable})}
Testés à chaque passe, dans cet ordre :
\begin{enumerate}
\item FC ventriculaire $< \vp{hr_arrest}$ (20/min) : \vp{cardiac_arrest} (aussi : moins de 3 battements après 10~s de conduction) ;
\item DC $< \vp{co_arrest_frac}\cdot 3{,}29\,(\vp{bw}/20)^{0{,}75}$ ($\vp{co_arrest_frac} = 0{,}15$) : \vp{circulatory_collapse} ;
\item PAM $< \vp{map_arrest}$ (40~mmHg) : \vp{irreversible_shock} (18/11/2026) ;
\item $\text{er} \ge 1$ (VO$_2$ > DO$_2$) : \vp{oxygen_debt} ;
\item échec d'une passe (racine introuvable, intégration impossible, FC ou DC non physiologiques) : \vp{no_steady_state} ;
\item pas de convergence en 30 passes : \vp{vicious_circle}.
\end{enumerate}

\subsection{Mécanismes de décompensation signalés}
L'état est « décompensé » si au moins un mécanisme est actif : \vp{acidemia_depression}, \vp{hypoxic_depression}, \vp{coronary_ischemia}, \vp{hypokalemic_weakness} (facteur $< 1$) ; \vp{atrial_standstill} ($K \ge 8{,}5$) ; \vp{conduction_block} (FC ventriculaire $< 0{,}98\times$ FC sinusale) ; \vp{anaerobic_metabolism} (er $> 0{,}55$) ; \vp{pulmonary_edema} (POG $> 18$) ; \vp{renal_hypoperfusion} (PAM $< 75$) ; \vp{baroreflex_exhausted} ($s_\text{tot} \ge 1$).

\begin{tabularx}{\linewidth}{@{}P{3.2cm}P{2.2cm}LP{0.7cm}@{}}
\caption{Paramètres de décompensation et de viabilité (\cmd{healthy\_dog.json}).}\label{tab:decomp}\\
\toprule
Paramètre & Valeur & Source & Év.\\
\midrule\endfirsthead
\toprule
Paramètre & Valeur & Source & Év.\\
\midrule\endhead
\vp{acid_ph0}, \vp{acid_ees}, \vp{acid_rsys} & 7,20 ; 1,5 ; 1,0 & aucune ; ordres de grandeur pédagogiques & D\\
\vp{hypox_pvo2}, \vp{hypox_ees} & 20 mmHg ; 0,03 & aucune & D\\
\vp{cpp_min}, \vp{cpp_ees} & 40 mmHg ; 0,02 & aucune & D\\
\vp{k_weak}, \vp{k_weak_vent} & 2,5 mmol/L ; 0,4 & aucune & D\\
\vp{hr_arrest}, \vp{co_arrest_frac} & 20 /min ; 0,15 & aucune & D\\
\vp{map_arrest} & 40 mmHg & PAM maintenue vers 40~mmHg : choc irréversible en 1 à 3~h (Morales 1999 ; Granot 1985 ; \acc{résumés}, sans DOI dans le dépôt) & C\\
\bottomrule
\end{tabularx}

\limite{Seuils et pentes inactifs dans la plage du golden et de la calibration, donc non contraints par les données du sain. Le critère de PAM est instantané alors que la source décrit une irréversibilité après 1 à 3~h. Le point fixe peut échouer près des frontières (cas « vicious\_circle » dépendant de la dose, section~\ref{sec:pk}). Le critère d'arrêt du baroréflexe ($|d_t - d| < 0{,}01$) laisse un écart résiduel de l'ordre de 0,5~mmHg entre deux calculs équivalents.}

%=====================================================================
\section{Signaux : ECG, VFC, auscultation}\label{sec:signaux}
%=====================================================================

\subsection{ECG (dérivation II)}
Somme de gaussiennes construite depuis la conduction (\cmd{signals::make\_ecg\_conducted}) : onde P à chaque activation atriale organisée (centre $t_a + 0{,}015$~s), ondes f en FA, QRS au temps de chaque activation ventriculaire, T après le QT corrigé. Pour un QRS fin de durée $q$ :
\begin{gather}
Q : -0{,}12\,R_\text{amp}\ \text{en}\ t - 0{,}34\,q,\ \sigma = q/14;\qquad R : R_\text{amp},\ \sigma = q/10;\qquad S : -0{,}20\,R_\text{amp}\ \text{en}\ t + 0{,}34\,q,\ \sigma = q/13,\\
T : T_\text{amp}\ \text{en}\ t - 0{,}5\,q + \text{QT}_c - 2{,}2\,t_w,\ \sigma = t_w,\qquad \text{QT}_c = \text{QT}\,(95/\max(60/RR, 30))^{1/3}\ \text{(Fridericia, RR précédent)}.
\end{gather}
P : $\sigma = \vp{P_SIGMA} = 10{,}5$~ms ; $t_w = \vp{T_W0} = 25$~ms. Onde Ta (repolarisation atriale, 18/11/2026) après chaque P organisée : amplitude $-0{,}35\,\vp{p_amp}$, $\sigma_{Ta} = 2{,}3\,\sigma_P$, centrée $2\sigma_P + \sigma_{Ta}$ après le centre de P (début à la fin de P, pic dans le segment PR) ; Perego 2014 (BAV~3, $n = 36$, \acc{résumé} : Ta/P $-0{,}35$ en amplitude, 2,3 en durée, $-0{,}09$~mV en D2) et Mirvis 1986 (pic dans le segment PR) ; masquée par le QRS en rythme sinusal (Romito 2022, \acc{lu}). Conduction intra-atriale \vp{pa_ia} = 41~ms (PA $41 \pm 1$~ms, AH $66 \pm 2$, HV $37 \pm 1$, $n = 97$ chiens anesthésiés ; Traunecker 1988, \acc{résumé}) : l'influx atteint le nœud AV $\vp{pa_ia}$ après le début de P et $\text{AH} = \text{PR} - \vp{hv} - \vp{pa_ia}$ (seul AH est modulé par le tonus autonome). Battement ventriculaire : QRS $2{,}2\,q$ (\vp{WIDE_QRS}), R $\times 1{,}2$ ($\sigma = 2{,}2q/8$), S $-0{,}35\,R_\text{amp}$ en $+0{,}40$, T discordante $-1{,}5\,T_\text{amp}$ de largeur $1{,}5\,t_w$. Onde f : amplitude $0{,}3\,\vp{p_amp}\cdot\min(2, \Delta t_a/\vp{af_cycle})$, $\sigma = 12$~ms. Bruit blanc 0,01~mV et dérive de ligne de base 0,02~mV à 0,25~Hz.

\textbf{Effets de K$^+$} (\cmd{ecg\_features}) : si $K > 5{,}5$ ($\delta = K - 5{,}5$) : $T_\text{amp} \times (1 + 0{,}45\delta)$ ; $t_w = \vp{T_W0}\max(0{,}025, 0{,}045 - 0{,}006\delta)/0{,}045$ ; P élargie $\times (1 + 0{,}25\max(0, K - 6))$ et aplatie $\times\max(0, 1 - 0{,}25\max(0, K - 6))$ ; QT $\times \max(0{,}8, 1 - 0{,}04\delta)$ ; PR $\times (1 + 0{,}07\max(0, K - 6))$ ; QRS $\times (1 + 0{,}12\max(0, K - 6{,}5))$. Si $K < 3{,}2$ : T $\times \max(0{,}2, 1 - 0{,}3(3{,}2 - K))$, QT $\times (1 + 0{,}08(3{,}2 - K))$. La disparition de P à 8,5~mmol/L vient de la conduction.

\sources{Largeurs calées pour que les durées visibles au seuil 0,02~mV soient les intervalles publiés : P 45~ms, QRS $= \vp{qrs}$, segment ST isoélectrique (Lerdweeraphon 2020, Vet World, doi:10.14202/vetworld.2020.1001-1004, \acc{lu} ; Bispo 2025, Vet Med Int, doi:10.1155/vmi/9947159, \acc{lu} ; Romito 2022, J Small Anim Pract, doi:10.1111/jsap.13532, \acc{lu} ; Santos 2024, \acc{lu}, intervalles cités de seconde main). Séquence hyperkaliémique (T ample, QT court, P plus longue et plus petite) : Pașca 2025, Vet Sci 12:63, doi:10.3390/vetsci12010063 (\acc{lu}, citation secondaire d'un traité). Critères des troubles du rythme : Kurogochi 2026 (FA), Santilli 2015 (BAV), Gunasekaran 2025 (ESV, pauses), Kanno 2017 (SSS), tous \acc{lu}. Placeholders D : \vp{T_W0}, \vp{P_WIDEN_K}, \vp{WIDE_QRS}, \vp{WIDE_R}, \vp{WIDE_T}, \vp{F_WAVE}, \vp{F_WAVE_MAX}, coefficients des effets de K$^+$.}
\limite{Dérivation II seule (pas d'axe électrique). Pas de perte de P ni de bloc à K 7,35 (cas publié : bloc complet à 7,9). Aucune donnée ECG d'hypokaliémie lue. Le QRS d'un battement ventriculaire fusionne avec sa T discordante sur le tracé.}

\subsection{VFC}
Voir la section~\ref{sec:auto} : les instants R viennent de la simulation autonome ; en rythme non sinusal, les signaux utilisent les décharges sinusales du modèle autonome, puis la conduction, puis l'ECG et les sons (\cmd{patient\_conduction}).

\subsection{Auscultation (PCG)}
\cmd{pcg.rs}, constantes toutes dans \cmd{modules/auscultation.json} (version 1.1.0), chacune avec \vp{source}, \vp{evidence} et \vp{calibre_sur}. Par battement :
\begin{itemize}
\item \textbf{B1} à la fermeture mitrale : trois sinusoïdes amorties ($f_0 = 70$~Hz $\times$ [0,8 ; 1,0 ; 1,45], poids [0,6 ; 1,0 ; 0,5], $\tau = 20$~ms $\times$ [1,0 ; 0,8 ; 0,6]), attaque de 4~ms, phases aléatoires ; amplitude $0{,}5\,\min(2,\ \text{d}P/\text{d}t_\text{max}/1800)\times$ gain de site ; seconde composante (tricuspidienne) retardée de 22~ms, amplitude relative 0,9.
\item \textbf{B2} aortique ($f_0$ 96~Hz) puis pulmonaire (90~Hz) retardée de $0{,}008 + 0{,}030\times$inspiration ; amplitude $0{,}4\,\min(2,\ (2P_{ao} - P_{vg})/190)$ ; systole $\times (RR/T)^{0{,}35}$, gain $\times (RR/T)^{0{,}5}$ ; garde-fou $t_{ac} \ge t_{mc} + 0{,}7\,t_s$.
\item \textbf{B3} si POG $> 15$~mmHg, 0,12~s après B2 : amplitude $0{,}35\max(0, (q_\text{précoce} - 160)/160)$ (\evD{}).
\item \textbf{Souffle} : bruit gaussien filtré (Butterworth 60--200~Hz, ordre 2), niveau $= 0{,}00054\,(q_\text{reg,max})^{1{,}15}$ si $q_\text{reg,max} > 3$~mL/s, forme $\sqrt{q_\text{reg}/q_\text{reg,max}}$, facteur par battement uniforme $[0{,}7; 1{,}2]$.
\item \textbf{Grade} indicatif : seuils $[10, 25, 50, 80, 120]$~mL/s sur le débit régurgitant maximal (\evD{}, non calibrés).
\item Variabilité lognormale ($\sigma$ 0,18) de B1 et B2, modulation respiratoire $1 + 0{,}15\sin$ ; bruit de fond 40--155~Hz (ordre 4, amplitude 0,014) modulé par la respiration. Sites : base (gains B1 0,7 ; B2 3,1) et apex (1,5 ; 3,1), souffle 0,7.
\end{itemize}
\sources{Heartsound Library de la TiHo Hannover (Vörös, Ehlers, Kleinsorgen, Nolte ; CC BY-NC-ND 4.0, aucun enregistrement embarqué) : 3 enregistrements sans souffle et 5 de régurgitation mitrale, un chien par grade. Calés (\evC{}) : durée de B1 (45 contre 43~ms), bruit de fond par bande, niveau du souffle par grade (écart quadratique moyen 3,4~dB sur 10 valeurs), rapport B2/B1 à l'apex (médiane 0,82). Systole : B1$\to$B2 du beagle (PMC9394128).}
\limite{B1$\to$B2 : 0,15~s depuis le recalage de la systole du 18/11/2026 (0,26~s avant), contre 0,18~s enregistré. B1 est placé au croisement P$_\text{VG}$/P$_\text{OG}$, que la systole atriale encore active au début de la contraction retarde (environ 60~ms après le début mécanique). Attribution du site « base » à revoir (l'enregistrement de calage est décrit comme pris à l'apex). Pas de souffle à la base ni de souffle diastolique enregistrés ; B3 non calé ; correspondance grade--débit non calibrée ; mp3 avec pertes.}

%=====================================================================
\section{Synthèse des sources et des placeholders}\label{sec:synthese}
%=====================================================================

{\small\begin{tabularx}{\linewidth}{@{}P{3.9cm}P{3.6cm}LP{1.3cm}@{}}
\caption{Sources utilisées par le modèle (accès déclaré dans le dépôt).}\label{tab:sources}\\
\toprule
Source & Identifiant & Utilisation & Accès\\
\midrule\endfirsthead
\toprule
Source & Identifiant & Utilisation & Accès\\
\midrule\endhead
Haskins 2005, Comp Med 55:156 & PMID 15884778 & FC, PAM, PAP, POG, IC, gaz du sain ; \vp{map_set} & résumé\\
Constable \& Stämpfli 2005 & \nolinkurl{10.1111/j.1939-1676.2005.tb02719.x} & Na, K, Cl, albumine, phosphate ; $A_\text{tot}$, p$K_a$ & résumé\\
Cambier 2004 & \nolinkurl{10.1016/j.rvsc.2004.02.004} & P50, effet Bohr & résumé\\
Bogucki \& Noszczyk-Nowak 2015 / 2017 & \nolinkurl{10.1515/pjvs-2015-0040} ; \nolinkurl{10.1515/pjvs-2017-0021} & VFC du sain (calage) ; MMVD C (comparaison) & lu\\
Billman 2013 & \nolinkurl{10.3389/fphys.2013.00222} & sens VFC--FC & lu\\
\nolinkurl{PMC9394128} & -- & \vp{ts_coef} (B1$\to$B2 du beagle) & n.p.\\
\nolinkurl{10.3390/ani15111524} & -- & \vp{k_la} (E/A) & n.p.\\
Schouten 1987 ; Wier 1986 ; Cooper 1993 ; Prinzen 1999 & voir section~\ref{sec:coeur} & relation force-intervalle, QRS large (structure) & n.p.\\
Yoran 1979 & non indiqué & orifice mitral dynamique (forme) & n.p.\\
Suzuki 2013 & \nolinkurl{10.1111/jvim.12154} (non revérifié) & POG sous furosémide (calage \vp{mr_k}, \vp{ven_cap}) & résumé\\
Bayne 1979 & non indiqué & veinodilatation du furosémide & n.p.\\
Koh 2021, BMC Vet Res 17:295 & \nolinkurl{10.1186/s12917-021-02998-4} & PK, diurèse, électrolytes du furosémide & lu\\
Wong 1977 & \nolinkurl{10.1007/BF03006233} (non vérifié) & K$^+$ sous furosémide répété & résumé\\
Miyazaki 1990 ; Hirai 1992 & \nolinkurl{10.1292/jvms1939.52.265} ; -- & effet piloté par l'excrétion urinaire & résumé\\
Adin 2020 ; Chetboul 2017 (TEST) & non indiqués & adaptation $W$ du furosémide & lu\\
Her 2020 & \nolinkurl{10.3389/fvets.2020.00423} & PK du pimobendane & lu\\
Mills 2025 ; Amon 2021 ; Wunsch 2010 & \nolinkurl{10.1111/jvp.13515} ; \nolinkurl{PMC7814539} ; \nolinkurl{PMC4235959} & PD de l'opioïde & lu\\
Maiante 2009 ; Garofalo 2012 ; Bailey 1987 ; Arndt 1984 ; Ingvast-Larsson 2010 & -- & PD, toxicité et PK de l'opioïde & résumé\\
Waterman 2019 ; Love 2000 & PMID 30801219 ; -- & profil aténolol & moteur\\
Collignon 2012 ; Bowsher 1985 & non indiqués & \vp{creat_vd}, \vp{urea_vd} & résumé\\
Morales 1999 ; Granot 1985 & non indiqués & \vp{map_arrest} & résumé\\
\nolinkurl{PMC7872775} ; Rigel 1990 & -- & \vp{av_symp} & lu ; résumé\\
Gunasekaran 2026 & \nolinkurl{10.3389/fcvm.2026.1890025} & \vp{av_erp}, \vp{av_conceal} ; module FA & n.p.\\
Gunasekaran 2025 & \nolinkurl{10.1371/journal.pone.0319886} & FC max des sains ; ESV ; pauses & lu\\
\nolinkurl{10.1111/jvim.12512} & -- & \vp{vent_escape_rate} 42/min & n.p.\\
Hariman 1983 ; Tag \& Day 2008 & -- & \vp{k_atrial_standstill} ; \vp{k_hp_block} & résumé ; moteur\\
ORCA & \nolinkurl{10.1111/jvim.16666} & FA non traitée & n.p.\\
\nolinkurl{10.3389/fvets.2025.1622850} & -- & seuil PR du BAV 1 & n.p.\\
\nolinkurl{10.1111/jvim.17071} ; \nolinkurl{10.1016/j.jvc.2018.12.003} & -- & pauses de la maladie du sinus & n.p.\\
Oyama 2004 ; Keene 2019 (ACVIM) ; MSD Vet Manual & \nolinkurl{10.1111/j.1939-1676.2004.tb02604.x} ; \nolinkurl{10.1111/jvim.15488} & MMVD : RF, stades, seuil d'œdème & résumé ; lu ; --\\
IRIS 2026 & iris-kidney.com & stades de la MRC & lu\\
Lerdweeraphon 2020 ; Bispo 2025 ; Romito 2022 ; Pașca 2025 & voir section~\ref{sec:signaux} & ECG & lu\\
TiHo Heartsound Library & -- & auscultation & mesures\\
\bottomrule
\end{tabularx}}
« n.p. » : niveau d'accès non précisé dans le dépôt. Les DOI des entrées de \cmd{references.bib} ont été vérifiés dans OpenAlex le 17/11/2026 ; les autres non, sauf mention.

\begin{tabularx}{\linewidth}{@{}P{3.4cm}LP{4.2cm}@{}}
\caption{Placeholders (evidence D) et valeurs sans source.}\label{tab:placeholders}\\
\toprule
Bloc & Valeurs & Remarque\\
\midrule\endfirsthead
\toprule
Bloc & Valeurs & Remarque\\
\midrule\endhead
Cœur & \vp{ts_restitution}, \vp{ca_tau}, \vp{ca_recirc}, \vp{ees_wide} ; \vp{ees_rv}, \vp{v0_*}, \vp{a_*}, \vp{b_rv} sans source & à caler sur d$P$/d$t$ vs couplage\\
Vaisseaux, volémie & compliances, volumes non contraints, \vp{r_valve}, volémie 85~mL/kg & calage global seulement\\
Oreillettes & \vp{k_ra}, \vp{ta_atr}, \vp{tra_atr} & \\
Baroréflexe & gains 0,8 / 0,6 / 0,4 / 0,5 / 0,20 / 45, \vp{a_s}, \vp{a_v}, \vp{s0}, \vp{v0}, \vp{baro_gain} & \\
VFC & constantes de temps, retards, sigmoïde, gain vasomoteur & calage global Bogucki 2015\\
Conduction & \vp{hv}, \vp{av_beta}, \vp{av_tau}, \vp{av_vagal}, \vp{av_k_gain}, \vp{af_cycle}, \vp{af_cv}, \vp{hp_k_slope}, \vp{junct_rate}, \vp{esc_k_gain}, \vp{esc_overdrive}, \vp{vpc_k_gain}, \vp{k_vpc} & \\
Gaz & contrôle ventilatoire, \vp{aagrad}, \vp{edema_k}, seuil 18 ; coefficient du phosphate (humain) ; \vp{hill_n} non vérifié & \\
Métabolisme & sigmoïde du lactate (0,55 ; 9 ; 0,06) & \\
Rein & fonction de perfusion, exposant 0,8 de l'urée ; \vp{gfr0}, \vp{creat0}, \vp{urea0} sans source & \\
Décompensation & tous les seuils sauf \vp{map_arrest} (C) & \\
PK/PD & PD de l'atropine, du bêtabloquant, du pimobendane ; PK de l'atropine et du bêtabloquant ; termes toxiques ; \vp{f_plasma}, \vp{tau_refill_h} ; EC$_{50}$ de l'opioïde & \\
Maladies & pneumonie (tout) ; ESV et maladie du sinus (coefficients) ; MRC hors créatinine ; écarts-types des tirages & \\
ECG & \vp{T_W0}, \vp{P_WIDEN_K}, \vp{WIDE_*}, \vp{F_WAVE*}, effets de K$^+$ & \\
PCG & B3, seuils de grade, garde-fou de systole, préréglage 72/min & \\
\bottomrule
\end{tabularx}



\section*{Points ouverts}
\addcontentsline{toc}{section}{Points ouverts}
\begin{itemize}
\item Les chiffres de sortie cités viennent des calages et des tests du 18/11/2026 (\cmd{tests/drug\_corrections.rs}, \cmd{docs/VERIFICATION\_MALADIES\_MEDICAMENTS.tex}) ; ils changent si un module est recalé.
\item Yoran 1979, Bayne 1979, Adin 2020, Chetboul 2017, Collignon 2012, Bowsher 1985, Morales 1999, Granot 1985 : cités sans DOI dans le dépôt ; ils ne figurent pas dans \cmd{references.bib}.
\item La provenance individuelle de nombreuses valeurs du chien sain (compliances, volumes non contraints, \vp{hb}, \vp{gfr0}, \vp{creat0}\ldots) n'est pas documentée ; le niveau de preuve par paramètre n'est pas un champ du JSON.
\end{itemize}

\end{document}
